% Auto-generated full Part-1 manuscript (pandoc + academic.latex).
% Source: docs/A_Unified_Theoretical_Framework_for_LENR.md (Part 1 only; Part 2 appendix excluded).
% Build: ./build_pdf.sh  (requires pandoc + XeLaTeX + xurl).
% Overflow fix: breakable \texttt via \allowbreak; wrapped AI contributor block; secnumdepth=-1.
% Epistemic: lenr_proved=false; all EQ-OPEN-* remain OPEN; no live ship / public flip.
%
\documentclass[11pt,a4paper]{article}

% Submission-oriented academic preamble (XeLaTeX; clean article class)
\usepackage[margin=1in]{geometry}
\usepackage{fontspec}
\setmainfont{Latin Modern Roman}
\setsansfont{Latin Modern Sans}
\setmonofont[Scale=0.90]{Latin Modern Mono}
\usepackage{amsmath,amssymb,amsfonts}
\usepackage{newunicodechar}
\newunicodechar{θ}{\ensuremath{\theta}}
\newunicodechar{γ}{\ensuremath{\gamma}}
\newunicodechar{δ}{\ensuremath{\delta}}
\newunicodechar{α}{\ensuremath{\alpha}}
\newunicodechar{β}{\ensuremath{\beta}}
\newunicodechar{π}{\ensuremath{\pi}}
\newunicodechar{σ}{\ensuremath{\sigma}}
\newunicodechar{ρ}{\ensuremath{\rho}}
\newunicodechar{ν}{\ensuremath{\nu}}
\newunicodechar{κ}{\ensuremath{\kappa}}
\newunicodechar{≠}{\ensuremath{\neq}}
\newunicodechar{∝}{\ensuremath{\propto}}
\newunicodechar{⁻}{\textsuperscript{-}}
\newunicodechar{⁰}{\textsuperscript{0}}
\newunicodechar{²}{\textsuperscript{2}}
\newunicodechar{⁴}{\textsuperscript{4}}
\newunicodechar{⁵}{\textsuperscript{5}}
\newunicodechar{₀}{\textsubscript{0}}
\newunicodechar{₂}{\textsubscript{2}}
\newunicodechar{∧}{\ensuremath{\wedge}}
\newunicodechar{↔}{\ensuremath{\leftrightarrow}}
\usepackage{booktabs}
\usepackage{graphicx}
\usepackage{xcolor}
\usepackage{setspace}
\usepackage{enumitem}
% xurl helps \url/\href; body paths use explicit \allowbreak via build_pdf.sh
\usepackage{xurl}
\urlstyle{tt}
\usepackage{hyperref}
\hypersetup{
  colorlinks=true,
  linkcolor=black,
  citecolor=black,
  urlcolor=blue,
  pdftitle={A Unified Theoretical and Mathematical Framework for LENR},
  pdfauthor={Diadon Acs}
}

% Manuscript headings already carry Part-1 numbers (1, 1.1, …).
\setcounter{secnumdepth}{-1}

% Prefer wrapping long identifiers/paths over margin overflow.
\tolerance=2500
\setlength{\emergencystretch}{4em}
\hbadness=10000


\setlength{\parskip}{0.35em}
\setlength{\parindent}{1.2em}
\setlist{nosep,leftmargin=*}
\onehalfspacing
\providecommand{\tightlist}{%
  \setlength{\itemsep}{0pt}\setlength{\parskip}{0pt}}

\title{A Unified Theoretical and Mathematical Framework for Low-Energy
Nuclear Reactions (LENR)\\\large\textit{Coupling is all you need}}
\author{Diadon Acs (lead, Conscious Energies)}
\date{2026-02-14 (revised 2026-11-11, EQ-OPEN peer-review refresh;~prior
2026-11-11 multidisciplinary coherence;~2026-11-10 Pyyhtiä CR-39 cite +
WS findings;~2026-11-05 honesty sync;~2026-09-25, v0.3.7-beta)}

\begin{document}
\maketitle

\begin{center}
\small Conscious Energies \textbar{} www.conscious.energy
\end{center}
\begin{center}
\footnotesize
Contributing Autonomous Intelligences (disclosed):\\[0.25em]
LENR-ARA/GPT-4, GPT-5.2, Sonnet 4.5, Sonnet 4.6,\\*
Opus 4.5, Opus 4.6, Opus 4.8, Gemini 2.5, Gemini 3.0,\\*
Grok 4.6, Composer (Cursor)\\[0.35em]
\textit{Human author solely accountable for content and claims.}
\end{center}
\vspace{0.5em}



\begin{abstract}
Low-Energy Nuclear Reactions (LENR) are investigated as a contested but
potentially consequential pathway toward clean and sustainable energy.
Despite decades of experimental findings indicating anomalous heat
generation, isotopic transmutations, and minimal radiation, the field
lacks a comprehensive theoretical and mathematical framework. This paper
aims to unify major theoretical elements---including Rydberg Matter
(RM), delocalized interface electrons, double-layer Quantum
Electrodynamics (QED) effects, Casimir effect, the Oppenheimer-Phillips
process, and quantum tunneling enhancements---under a single model.
Lattice-bubble interaction physics (heterogeneous nucleation at
surfaces, bubble--surface coupling, thermal and diffusion boundary
layers, and acoustic/ultrasonic forcing) is integrated as a
\textbf{modelled enhancement / contested hydro} candidate environment
for cavitation-assisted interfacial conditions---not as proof of nuclear
rates, and not as a claim that laboratory cavitation establishes LENR.
In tandem, we detail mathematical tools and simulations that quantify
these LENR hypotheses under explicit claim tiers. Computational
explorations in the companion LENR\_SIM software implement a
phenomenological tunneling-enhancement band of \(10^{3}\)--\(10^{5}\)
(revised downward from earlier \(10^{8}\) estimates; see Section 5.1);
this band is a software sensitivity envelope, not a first-principles
prediction. Reported isotopic shifts, neutron signals, and excess heat
remain disputed empirical correlations that motivate---but do not
validate---the framework. A companion Lean 4 formalization (no mathlib;
no \texttt{sorry}; one labelled RealShim axiom) records identities, DEF
clamps, and negative/bound results under stated assumptions; it does
\textbf{not} prove LENR ({\ttfamily lenr\_\allowbreak{}proved\ =\allowbreak{}\ false}). All five
\texttt{EQ-OPEN-*} obligations remain open: affirmative closure requires
governed experimental evidence with calibration and provenance, not Lean
alone.
\end{abstract}

\noindent\textit{Keywords:} LENR; electron--neutron--proton
coupling; Rydberg Matter; quantum tunneling; Casimir; claim tiers.

\vspace{0.75em}
\noindent\footnotesize
\textbf{Honesty footnotes.}
\texttt{lenr\_proved} remains false.
All five \texttt{EQ-OPEN-*} obligations remain OPEN
(\texttt{EQ-OPEN-MEV-01}, \texttt{EQ-OPEN-RAD-01}, \texttt{EQ-OPEN-RATE-01},
\texttt{EQ-OPEN-EC-01}, \texttt{EQ-OPEN-YIELD-01}).
Cavitation--collapse is modelled enhancement / contested hydro only.
Quantitative bands unchanged: tunneling $10^{3}$--$10^{5}$;
Casimir $-1.3\,\mathrm{mN}$; Miles $24\pm4\,\mathrm{MeV}$.
No live ship / public flip is implied by this archival PDF.
\normalsize

\begin{quote}
\textbf{Epistemic posture (submission).} Exploratory models only.
Claim-tiered outputs. \texttt{lenr\_proved} remains \textbf{false}. All
five \texttt{EQ-OPEN-*} obligations (\texttt{EQ-OPEN-MEV-01},
\texttt{EQ-OPEN-RAD-01}, \texttt{EQ-OPEN-RATE-01},
\texttt{EQ-OPEN-EC-01}, \texttt{EQ-OPEN-YIELD-01}) remain \textbf{OPEN}.
Cavitation--collapse hydrodynamics is \textbf{modelled enhancement /
contested hydro}, not nuclear-rate proof. No quantitative band changes
(tunneling \(10^{3}\)--\(10^{5}\); Casimir \(-1.3\,\mathrm{mN}\); Miles
\(24\pm4\,\mathrm{MeV}\)).
\end{quote}

\hypertarget{introduction}{%
\subsection{1. Introduction}\label{introduction}}

\hypertarget{background-and-challenges}{%
\subsubsection{1.1 Background and
Challenges}\label{background-and-challenges}}

Low-Energy Nuclear Reactions (LENR) refer to nuclear processes
hypothesized to occur at or near room temperature and pressure,
challenging the conventional high-energy paradigm of nuclear fusion or
fission. Over the past several decades, experiments have reported:

\begin{itemize}
\tightlist
\item
  \textbf{Anomalous Heat Generation:} Measured energy outputs exceeding
  what chemical processes alone can produce {[}1,4{]}.\\
\item
  \textbf{Isotopic Transmutations:} Conversion of elements without the
  high-energy nuclear by-products usually observed {[}2,5{]}.\\
\item
  \textbf{Minimal Radiation:} Comparatively negligible ionizing
  radiation despite \emph{reported} nuclear-scale heat / isotopic
  correlations {[}3,6{]}---reports that remain disputed and do not by
  themselves establish nuclear events (\texttt{EQ-OPEN-RAD-01} stays
  OPEN).
\end{itemize}

However, the field faces two major hurdles:

\begin{itemize}
\tightlist
\item
  \textbf{Reproducibility:} Varying experimental conditions lead to
  inconsistent results across different laboratories {[}3,15{]}.\\
\item
  \textbf{Theoretical Fragmentation:} Absence of a unifying theoretical
  and mathematical framework that convincingly explains the diverse
  observed phenomena {[}18,22{]}.
\end{itemize}

Given LENR's contested energy relevance, resolving these issues remains
a priority for scientific understanding. Energetic / practical freedom
outcomes stay a permanent non-goal of the companion software governance
(see \texttt{GOALS.md}).

\hypertarget{objectives}{%
\subsubsection{1.2 Objectives}\label{objectives}}

This paper aims to:

\begin{enumerate}
\def\labelenumi{\arabic{enumi}.}
\tightlist
\item
  \textbf{Unify Diverse Phenomena:} Integrate Rydberg Matter (RM), halo
  electrons, double-layer QED, Casimir forces, the Oppenheimer-Phillips
  process, and quantum tunneling enhancements into a single LENR
  \textbf{hypothesis stack} under explicit claim tiers.\\
\item
  \textbf{Quantify Mechanisms:} Use detailed mathematical formulations,
  computational simulations, and statistical analysis to quantify each
  key mechanism as \textbf{modelled} sensitivity exploration unless an
  independent experimental constraint is cited.\\
\item
  \textbf{Propose Reproducible Protocols:} Offer experimental procedures
  and theoretical guidelines that improve laboratory reproducibility and
  independent falsification---not industrial or commercial LENR
  applications.
\end{enumerate}

\begin{center}\rule{0.5\linewidth}{0.5pt}\end{center}

\hypertarget{theoretical-framework}{%
\subsection{2. Theoretical Framework}\label{theoretical-framework}}

\hypertarget{rydberg-matter-rm}{%
\subsubsection{2.1 Rydberg Matter (RM)}\label{rydberg-matter-rm}}

\textbf{Claim tier:} speculative / exploratory mechanism (lattice--RM
hypothesis stack). Lean \texttt{LenrFormal.Rydberg} locks the \(n=10\)
energy identity only---not a nature proof that RM mediates LENR.

\textbf{Definition:} Rydberg Matter is a phase in which atoms or
molecules occupy highly excited electronic states with extended orbitals
{[}21{]}.

\begin{enumerate}
\def\labelenumi{\arabic{enumi}.}
\item
  \textbf{Role in LENR}

  \begin{itemize}
  \tightlist
  \item
    \textbf{Enhanced Energy Localization:} Extended orbitals allow more
    efficient concentration of energy, \emph{hypothesized} to lower the
    Coulomb barrier for nuclear processes.\\
  \item
    \textbf{Formation Sites:} Commonly proposed at metal--hydrogen
    interfaces during electrolysis or near catalytic surfaces in gaseous
    systems under specific pressures and temperatures.\\
  \item
  \end{itemize}
\item
  \textbf{Energy Calculation}\\
  For an atom at principal quantum number \(n = 10\):

  \[
   E_n = -\frac{13.6\,\mathrm{eV}}{n^2} = -0.136\,\mathrm{eV},
   \]

  indicating a weaker binding energy than ground-state levels, which can
  promote higher electron mobility.
\end{enumerate}

\begin{center}\rule{0.5\linewidth}{0.5pt}\end{center}

\hypertarget{delocalized-interface-electrons-extended-electronic-states}{%
\subsubsection{2.2 Delocalized Interface Electrons (Extended Electronic
States)}\label{delocalized-interface-electrons-extended-electronic-states}}

\textbf{Claim tier:} modelled / exploratory construct. Lean
\texttt{LenrFormal.Overlap} records ensemble-count bookkeeping
(\texttt{EQ-OVERLAP-01/02}); ensemble product ≠ nuclear rate.

\textbf{Definition:} In metallic hydrides and at metal-hydrogen
interfaces, electrons can occupy extended or delocalized quantum states
with wavefunctions that have non-negligible amplitude at nuclear length
scales. These ``extended electronic states'' arise from band-structure
effects, surface states, and hybridization in the condensed matter
environment. (Citation correction: an earlier draft cited reference
{[}12{]} (Meng et al.~2013) here, but that work concerns \emph{halo
nuclei}---nuclei with weakly bound nucleons extending beyond the nuclear
core---and does \textbf{not} support a claim about electrons. Reference
{[}12{]} is therefore \emph{not} cited as evidence for delocalized
interface electrons; this concept rests instead on standard
condensed-matter band-structure and surface-state physics, and is
treated here as a model construct whose overlap magnitude is bounded in
Section 5.1 rather than asserted from {[}12{]}.)

\begin{enumerate}
\def\labelenumi{\arabic{enumi}.}
\item
  \textbf{Overlap Probability}\\
  The probability of an electronic wavefunction overlapping with the
  nuclear volume is:

  \[
   P_{\text{overlap}} = \int_{0}^{r_0} |\psi_e(r)|^2 \, 4\pi r^2 \, dr
   \]

  For extended electronic states with characteristic orbital scale
  \(a = 10\,a_0 = 0.529\) nm and the nuclear interaction radius
  \(r_0 \approx 5\) fm \(= 5 \times 10^{-6}\) nm (the nuclear radius of
  deuterium, \(r_0 \approx r_d \approx 1.2 A^{1/3}\) fm \(\approx 1.5\)
  fm, though we use \(\sim 5\) fm as an effective nuclear interaction
  range), the overlap probability is approximated by:

  \[
   P_{\text{overlap}} \approx \left(\frac{r_0}{a}\right)^3 e^{-2r_0/a}
   \]

  Since \(r_0/a = 5 \times 10^{-6} / 0.529 = 9.45 \times 10^{-6}\), the
  exponential factor is
  \(e^{-2 \times 9.45 \times 10^{-6}} \approx 1.0\), giving:

  \[
   P_{\text{overlap}} \approx (9.45 \times 10^{-6})^3 \approx 8.5 \times 10^{-16}
   \]

  This is consistent with the Part 2 LaTeX analysis of this paper, which
  uses \(\sim 1\) fm scales and obtains \(P \sim 10^{-10}\) to
  \(10^{-12}\) per event (the difference reflecting different
  assumptions about the effective wavefunction amplitude at the
  nucleus).

  \textbf{Ensemble scaling (order-of-magnitude only).} Individually,
  \(P_{\text{overlap}} \sim 10^{-16}\) (or \(10^{-10}\)--\(10^{-12}\)
  under alternate nuclear-radius assumptions) is negligible. Multiplying
  by macroscopic atom counts (\(\sim 10^{23}\)) and high attempt
  frequencies does \textbf{not} by itself establish nuclear reaction
  rates: most ``overlap events'' are not reaction events, and the
  conversion factor is precisely the unknown matrix element under test.
  The following product is therefore an \textbf{upper-bound counting
  argument}, not a predicted LENR rate:

  \[
   R_{\text{overlap}} \sim 10^{23} \times 10^{12} \times 10^{-12} \sim 10^{23} \text{ potential overlap events per second}
   \]

  This large statistical ensemble is a key feature distinguishing
  condensed matter nuclear processes from few-body scattering
  experiments.
\item
  \textbf{Enhanced Nuclear Interaction}\\
  This non-zero overlap at nuclear scales enables weak-interaction
  processes such as electron capture (\(p + e^- \to n + \nu_e\)), which
  is an essential element in proposed LENR transmutation pathways. The
  rate depends critically on the electron density at the nucleus,
  \(|\psi_e(0)|^2\), which can be enhanced by lattice confinement,
  coherent many-body effects, and the anomalous screening discussed in
  Section 2.1.
\end{enumerate}

\begin{center}\rule{0.5\linewidth}{0.5pt}\end{center}

\hypertarget{double-layer-qed-effects}{%
\subsubsection{2.3 Double-Layer QED
Effects}\label{double-layer-qed-effects}}

\textbf{Claim tier:} speculative / exploratory QED--lattice coupling.
Lean \texttt{LenrFormal.Field} / \texttt{ScaleQED} document identities
and negative controls; bulk HE vacuum polarisation remains a negative
control (\texttt{EQ-QED-HE-NEG-01}), not LENR proof.

\textbf{Definition:} Charge separation at electrolyte or lattice
interfaces can form high electric fields in thin double layers
{[}23,33{]}.

\begin{itemize}
\tightlist
\item
  \textbf{Field Calculation:} For a potential difference of
  \(\Delta V = 0.1\,\text{V}\) over a distance \(d = 1\,\text{nm}\):
\end{itemize}

\[
E = \frac{\Delta V}{d} = \frac{0.1\,\text{V}}{1\times 10^{-9}\,\text{m}} = 1\times 10^{8}\,\text{V/m}.
\]

\begin{itemize}
\item
  \textbf{Implications:}

  \begin{itemize}
  \tightlist
  \item
    Coulomb Barrier Reduction through localized high fields.\\
  \item
    Ionization enhancement and potential formation of Rydberg Matter.
  \end{itemize}
\end{itemize}

\hypertarget{quantum-field-theory-foundation}{%
\paragraph{2.3.1 Quantum Field Theory
Foundation}\label{quantum-field-theory-foundation}}

The double-layer QED effects can be rigorously formulated using quantum
electrodynamics in condensed matter. The fundamental QED Lagrangian
density is:

\[
\mathcal{L} = \bar{\psi}(i\gamma^\mu D_\mu - m)\psi - \frac{1}{4}F_{\mu\nu}F^{\mu\nu} + \mathcal{L}_{\text{int}}
\]

where: - \(\bar{\psi}, \psi\) are the electron field operators -
\(\gamma^\mu\) are the Dirac matrices -
\(D_\mu = \partial_\mu + ieA_\mu\) is the covariant derivative -
\(F_{\mu\nu} = \partial_\mu A_\nu - \partial_\nu A_\mu\) is the
electromagnetic field tensor - \(\mathcal{L}_{\text{int}}\) represents
additional interactions (lattice, phonon coupling)

\textbf{Vacuum Fluctuations and Casimir Connection:}

The vacuum energy density in the presence of boundaries (e.g., at the
double layer interface) is:

\[
\langle 0 | T^{00} | 0 \rangle = \frac{1}{2} \sum_k \hbar \omega_k
\]

In confined geometries, the mode spectrum is altered, leading to the
Casimir effect. The vacuum polarization in strong fields
(\(E \sim 10^8\) V/m) modifies the effective dielectric constant:

\[
\epsilon_{\text{eff}} = 1 + \frac{2\alpha}{45\pi} \left(\frac{E}{E_{\text{cr}}}\right)^2
\]

where \(\alpha \approx 1/137\) is the fine structure constant and
\(E_{\text{cr}} = m_e^2 c^3 / (e\hbar) \approx 1.3 \times 10^{18}\) V/m
is the critical field.

\textbf{Field-Matter Coupling in Condensed Systems:}

In metallic hydrides, the electron-phonon-photon coupling can be
expressed as:

\[
\mathcal{L}_{\text{int}} = g_{e-ph} \bar{\psi} \psi \phi + g_{ph-\gamma} \phi F_{\mu\nu} F^{\mu\nu}
\]

where \(\phi\) represents the phonon field, and \(g_{e-ph}\),
\(g_{ph-\gamma}\) are coupling constants. The second term couples the
phonon field to the electromagnetic field strength tensor
\(F_{\mu\nu} = \partial_\mu A_\nu - \partial_\nu A_\mu\), ensuring gauge
invariance of the interaction. This formulation provides a rigorous
foundation for understanding how electromagnetic fields in double layers
couple to lattice dynamics and ultimately influence nuclear processes.

\textbf{Implications for LENR:}

The QED framework predicts that: 1. Strong local fields can polarize the
vacuum and modify tunneling barriers 2. Coherent photon modes can couple
to lattice phonons, enabling energy transfer 3. The collective
electromagnetic response of electrons in the double layer can enhance
screening effects beyond simple electrostatic models

\begin{center}\rule{0.5\linewidth}{0.5pt}\end{center}

\hypertarget{neutron-formation-via-electron-capture}{%
\subsubsection{2.4 Neutron Formation via Electron
Capture}\label{neutron-formation-via-electron-capture}}

The reaction:

\[
p + e^- \longrightarrow n + \nu_e
\]

is a \textbf{weak interaction} process (inverse beta decay) with a
threshold energy of:

\[
\Delta E = (m_n - m_p - m_e)c^2 \approx 0.782 \text{ MeV}
\]

This threshold means free electrons at thermal energies cannot drive
this reaction; the electron must have kinetic energy exceeding 0.782
MeV, or the energy must be supplied by other means (e.g., lattice
momentum, coherent excitation, or effective mass renormalization).

\begin{itemize}
\item
  \textbf{Cross-Section for Inverse Beta Decay.} Near threshold, the
  cross-section for this weak-interaction process is given by:

  \[
    \sigma \sim \frac{G_F^2}{\pi} \frac{E_\nu^2}{(\hbar c)^4}
    \]

  where \(G_F \approx 1.166 \times 10^{-5}\) GeV\(^{-2}\) is the Fermi
  coupling constant and \(E_\nu\) is the outgoing neutrino energy. For
  electrons with energy near threshold (\(\sim 1\) MeV), this yields:

  \[
    \sigma \sim 10^{-47} \text{ m}^2
    \]

  This is \textbf{26 orders of magnitude smaller} than the previously
  stated value of \(10^{-21}\) m\(^2\), which corresponds to
  electromagnetic or strong interaction scales. Weak interaction
  cross-sections are characteristically tiny, which is precisely why
  neutron formation via electron capture is ordinarily negligible.
\item
  \textbf{Widom-Larsen ``Heavy Electron'' Hypothesis.} Widom and Larsen
  (2006) proposed that collective electromagnetic excitations in metal
  hydride surfaces could renormalize the electron mass to an effective
  value \(m_e^* \gg m_e\), effectively lowering the kinematic threshold
  for this reaction. For the reaction to become energetically allowed at
  rest, one requires:

  \[
    m_e^* c^2 > (m_n - m_p)c^2 = 1.293 \text{ MeV} \quad \Rightarrow \quad \frac{m_e^*}{m_e} > \frac{1.293}{0.511} \approx 2.53
    \]

  This mechanism remains controversial: while the electromagnetic field
  energy near surfaces can be large, whether it coherently renormalizes
  the electron mass to this degree in condensed matter systems is
  debated
  \[Hagelstein et al., J. Condensed Matter Nucl. Sci. 12, 105 (2013)\].
\item
  \textbf{Corrected Rate Estimate.} The neutron production rate per unit
  volume is:

  \[
    R = n_p \, n_e \, \langle \sigma v \rangle
    \]

  where \(n_p\) and \(n_e\) are the proton and electron number densities
  and the thermal average \(\langle \sigma v \rangle\) accounts for the
  energy threshold. For standard thermal electrons, the Boltzmann
  suppression factor
  \(\exp(-\Delta E / k_B T) \sim \exp(-0.782 \text{ MeV} / 0.026 \text{ eV}) \sim e^{-3 \times 10^{4}}\)
  renders this rate effectively zero, underscoring that some
  non-equilibrium or collective mechanism is required if this process is
  to contribute meaningfully to LENR.
\item
  \textbf{Implications:} The extreme smallness of the weak-interaction
  cross-section and the energy threshold make neutron formation via
  simple electron capture negligible under ordinary conditions. Any
  viable LENR model invoking this pathway must explain how the effective
  threshold is overcome---whether through coherent many-body effects,
  mass renormalization, or other non-equilibrium mechanisms. Companion
  Lean \texttt{LenrFormal.WeakInteraction} / IR
  {\ttfamily EQ-\allowbreak{}EC-\allowbreak{}THRESH/\allowbreak{}GF/\allowbreak{}SIGMA/\allowbreak{}RATE-\allowbreak{}01} and ledger \texttt{ENG-WEAK-01}
  lock PDG \(G_F\), a \(\sigma\lesssim 10^{-47}\,\mathrm{m}^{2}\)
  envelope, and a Boltzmann-exponent floor as \textbf{identities /
  bounds only}---they do \textbf{not} discharge \texttt{EQ-OPEN-EC-01},
  assert a condensed-matter EC path, or flip \texttt{lenr\_proved}.
\end{itemize}

\begin{center}\rule{0.5\linewidth}{0.5pt}\end{center}

\hypertarget{casimir-effect}{%
\subsubsection{2.5 Casimir Effect}\label{casimir-effect}}

\textbf{Definition:} Quantum vacuum fluctuations produce an attractive
force between closely spaced conductive or dielectric surfaces {[}17{]}.

\begin{itemize}
\item
  \textbf{Force Calculation:} For area \(A = 1\,\mu\text{m}^2\) and
  separation \(d = 1\,\text{nm}\): \[
   F_{\text{Casimir}} = -\frac{\pi^2 \hbar c}{240}\,\frac{A}{d^4}
   \]

  Using \(\hbar = 1.055 \times 10^{-34}\) J·s, \(c = 3.00 \times 10^8\)
  m/s, \(A = 1 \times 10^{-12}\) m², and \(d = 1 \times 10^{-9}\) m:

  \[
   F_{\text{Casimir}} = -\frac{9.87 \times 1.055 \times 10^{-34} \times 3.00 \times 10^8}{240} \times \frac{10^{-12}}{(10^{-9})^4} \approx -1.3 \times 10^{-3} \text{ N}
   \]

  \textbf{Corrected Note:} The value is approximately \(-1.3\) mN (not
  \(-10^{-7}\) N). For \(-10^{-7}\) N, either
  \(A = 7.7 \times 10^{-16}\) m² or \(d = 3.16\) nm would be required.\\
  (\emph{Exact values vary by geometry and material; this is an
  order-of-magnitude estimate.})
\item
  \textbf{Implications:}

  \begin{itemize}
  \tightlist
  \item
    \textbf{Energy Localization:} Casimir forces can induce nanoscale
    distortions, focusing energy that can facilitate nuclear events.\\
  \item
    \textbf{Enhanced Tunneling:} Confinement at the nanoscale further
    lowers effective barriers.
  \end{itemize}
\end{itemize}

\hypertarget{quantum-coherence-length-scale}{%
\paragraph{2.5.1 Quantum Coherence Length
Scale}\label{quantum-coherence-length-scale}}

For coherent tunneling clusters, the coherence length determines the
spatial extent over which quantum correlations persist. This is a
critical parameter for understanding multi-body coherent effects in
LENR.

\textbf{Electronic coherence length.} For electrons mediating screening
and collective effects, the coherence length is set by the electron
Fermi velocity:

\[
\xi_e = \frac{\hbar v_F}{\Delta E}
\]

where \(v_F\) is the electron Fermi velocity and \(\Delta E\) is the
relevant energy gap. For LENR conditions in transition metals:

\begin{itemize}
\tightlist
\item
  \(v_F \approx 10^6\) m/s (typical for Pd, Ni)
\item
  \(\Delta E \approx 0.1\text{-}1\) eV (energy scale of electronic
  processes)
\end{itemize}

\[
\xi_e \approx \frac{1.055 \times 10^{-34} \times 10^6}{0.1 \times 1.602 \times 10^{-19}} \approx 6.6 \text{ nm}
\]

This electronic coherence length sets the spatial scale over which
electron-mediated screening and collective electromagnetic effects can
persist.

\textbf{Nuclear (deuteron) coherence length.} For deuteron-deuteron
quantum coherence---the directly relevant quantity for coherent
tunneling and nuclear fusion---the appropriate velocity is the deuteron
thermal velocity, not the electron Fermi velocity:

\[
v_D = \sqrt{\frac{k_B T}{m_D}} \approx \sqrt{\frac{1.381 \times 10^{-23} \times 300}{3.344 \times 10^{-27}}} \approx 1{,}100 \text{ m/s}
\]

\[
\xi_D = \frac{\hbar v_D}{\Delta E} \approx \frac{1.055 \times 10^{-34} \times 1100}{0.1 \times 1.602 \times 10^{-19}} \approx 7.2 \text{ pm}
\]

This sub-angstrom coherence length implies that deuteron-deuteron
quantum coherence is confined to nearest-neighbor pairs at thermal
equilibrium. For coherent multi-body fusion (e.g., Takahashi's TSC
model) to operate over nanometer scales, an additional mechanism is
required to extend nuclear coherence---for example, the non-equilibrium
BEC framework discussed in Section 2.6, which could dramatically
increase the deuteron de Broglie wavelength through effective cooling
and mass renormalization.

\textbf{Implication for simulation:} The coherent cluster size in the
simulation must be constrained by the appropriate coherence length. For
electron-mediated effects, clusters up to \textasciitilde6.6 nm are
physically justified. For direct nuclear coherence, cluster sizes must
be limited to \(\xi_D\) unless the non-equilibrium BEC conditions of
Section 2.6 are satisfied, in which case the effective coherence length
may extend to the electronic scale.

\begin{center}\rule{0.5\linewidth}{0.5pt}\end{center}

\hypertarget{bose-einstein-condensation-bec}{%
\subsubsection{2.6 Bose-Einstein Condensation
(BEC)}\label{bose-einstein-condensation-bec}}

\textbf{Definition:} BEC occurs when particles coalesce into a single
quantum state, enabling coherent interactions {[}25{]}.

\begin{itemize}
\item
  \textbf{Critical Temperature Estimate:} For a density
  \(n = 1 \times 10^{28}\) m⁻³ of deuterons (mass
  \(m = 3.344 \times 10^{-27}\) kg):

  \[
   T_c = \frac{2\pi \hbar^2}{m k_B} \left(\frac{n}{\zeta(3/2)}\right)^{2/3}
   \]

  Using:

  \begin{itemize}
  \tightlist
  \item
    \(\hbar = 1.055 \times 10^{-34}\) J·s
  \item
    \(m = 3.344 \times 10^{-27}\) kg (deuteron mass)
  \item
    \(k_B = 1.381 \times 10^{-23}\) J/K
  \item
    \(\zeta(3/2) \approx 2.612\)
  \end{itemize}

  \[
   T_c = \frac{2\pi \times (1.055 \times 10^{-34})^2}{3.344 \times 10^{-27} \times 1.381 \times 10^{-23}} \times \left(\frac{10^{28}}{2.612}\right)^{2/3} \approx 7.9 \text{ K}
   \]
\end{itemize}

\textbf{Critical Assessment and Non-Equilibrium Extension:}

The equilibrium BEC temperature of 7.9 K is far below LENR operating
conditions (300-800 K), making equilibrium BEC impossible. However,
non-equilibrium driven-dissipative mechanisms may achieve the required
phase-space density through three synergistic effects:

\begin{enumerate}
\def\labelenumi{\arabic{enumi}.}
\item
  \textbf{Effective Mass Renormalization (Polaron Dressing):} Deuterons
  strongly coupled to the phonon cloud form polaronic quasiparticles
  with effective mass \(m^* < m_D\). For strong electron-phonon coupling
  in PdD systems, mass reduction factors of 3-5x are plausible, raising
  \(T_c\) to 24-40 K. \emph{(Model assumption: the 3-5x factor is an
  order-of-magnitude estimate, not a measured value; it is treated as an
  uncertain input for the Monte Carlo / sensitivity analysis of Section
  5.1 and is not independently cited.)}
\item
  \textbf{Driven Dissipative Cooling (Phonon Standing Wave Mechanism):}
  Coherent phonon standing waves---whether acoustically driven or
  arising from lattice resonances---create a spatially periodic
  potential for interstitial deuterons. Through anharmonic coupling,
  deuteron kinetic energy is preferentially channeled into
  high-frequency optical phonon modes that rapidly dissipate into the
  bulk lattice. This mechanism is analogous to resolved sideband cooling
  in optomechanics and Sisyphus cooling in atomic physics. Estimated
  effective temperature reduction: factor of 3-5x, giving
  \(T_{\text{eff}} \sim 60\text{-}100\) K.
\item
  \textbf{Local Density Enhancement at Standing Wave Nodes:} Phonon
  radiation pressure concentrates deuterons at nodes or antinodes of the
  standing wave pattern, locally increasing number density by factors of
  5-10x. Since \(T_c \propto n^{2/3}\), this raises \(T_c\) by factors
  of 2.9-4.6x.
\end{enumerate}

\textbf{Combined Quantitative Estimate:}

\[
T_c^{\text{eff}} \approx 7.9 \text{ K} \times 5 \text{ (mass)} \times 4.6 \text{ (density)} \approx 182 \text{ K}
\]

\[
T_{\text{deuteron}}^{\text{eff}} \approx 300 \text{ K} / 5 \text{ (cooling)} \approx 60 \text{ K}
\]

Since \(T_{\text{deuteron}}^{\text{eff}} < T_c^{\text{eff}}\), BEC
becomes formally possible under combined non-equilibrium conditions.

\textbf{Open Constraints:} (a) Phonon-phonon thermalization in metals
occurs on \textasciitilde ps timescales; the cooling mechanism must
cycle faster than thermalization re-equilibrates. (b) The direct
criterion is phase-space density
\(n \lambda_{\text{dB}}^3 \geq \zeta(3/2)\) rather than temperature
comparison alone. (c) Whether the resulting state is true BEC
(spontaneous symmetry breaking) or a coherently driven state (analogous
to a laser) affects the predicted nuclear enhancement and must be
distinguished experimentally.

\begin{center}\rule{0.5\linewidth}{0.5pt}\end{center}

\hypertarget{oppenheimer-phillips-process}{%
\subsubsection{2.7 Oppenheimer-Phillips
Process}\label{oppenheimer-phillips-process}}

\textbf{Definition:} A mechanism where a deuteron interacts with a
nucleus at low energies, resulting in neutron capture and proton
emission {[}16{]}.

\[
d + {}^{A}_{Z}\text{X} \longrightarrow {}^{A+1}_{Z}\text{Y} + p.
\]

\begin{enumerate}
\def\labelenumi{\arabic{enumi}.}
\item
  \textbf{Mechanism}

  \begin{itemize}
  \tightlist
  \item
    Neutron from the deuteron is absorbed, forming
    \({}^{A+1}_{Z}\text{Y}\).\\
  \item
    The proton is repelled by Coulomb forces and emitted.\\
  \end{itemize}
\item
  \textbf{Rate Calculation}
\end{enumerate}

\[
R_{\text{rate}} = n_d\,n_X\,\langle\sigma_{\text{OP}} v_d\rangle,
\]

where \(\sigma_{\text{OP}}\) is the Oppenheimer-Phillips
(deuteron-stripping) cross-section, \(n_d\) and \(n_X\) are deuteron and
target densities, and \(\langle\sigma_{\text{OP}} v_d\rangle\) is the
reactivity formed with the deuteron velocity distribution (not a
free-neutron flux).

\begin{center}\rule{0.5\linewidth}{0.5pt}\end{center}

\hypertarget{quantum-tunneling-enhancements}{%
\subsubsection{2.8 Quantum Tunneling
Enhancements}\label{quantum-tunneling-enhancements}}

\textbf{Standard Tunneling Probability}

\[
P_{\text{tunnel}} = \exp\left(-\frac{2}{\hbar}\int_{a}^{b}\sqrt{2m\bigl(V(x)-E\bigr)}\,dx\right).
\]

\textbf{LENR Conditions}

\begin{itemize}
\tightlist
\item
  \textbf{Reduced Barrier Height:} Electric fields, Rydberg states, and
  Casimir effects lower \(V(x)\).\\
\item
  \textbf{Enhanced Particle Energy:} Localized heating and lattice
  phonon interactions can further increase effective particle energies.
\end{itemize}

\textbf{Overall Enhancement}\\
The reported \(10^{3}\)--\(10^{5}\) range is a \textbf{bounded
phenomenological sensitivity band} implemented in the current LENR\_SIM
software (independent mechanism multipliers with a hard composition
cap), \textbf{not} a first-principles prediction of a coupled
electron--lattice--nuclear Hamiltonian. Factorization into electron
screening (\(f_{\text{screening}} \sim 10^{1-2}\)), interface electric
fields (\(f_{\text{field}} \sim 10^{0.5-1}\)), lattice coherence
(\(f_{\text{coherence}} \sim 10^{1-2}\)), and phonon coupling
(\(f_{\text{phonon}} \sim 10^{0-0.5}\)) is a bookkeeping device:
\(f_{\text{total}} = f_{\text{screening}} \times f_{\text{field}} \times f_{\text{coherence}} \times f_{\text{phonon}} \sim 10^{3} \text{ to } 10^{5}\),
subject to software caps and claim-tier labels.\footnote{Software
  composition is \textbf{not} this four-factor product. The theory
  pipeline uses grouped-max then a \(10^{5}\) cap (\texttt{EQ-COMP-01},
  \texttt{EQ-COMP-02}); multiplying that total by the bubble-lattice
  peak (or the legacy geometric-mean path) is an illegal product
  (\texttt{EQ-COMP-03}). See
  {\ttfamily LENR\_\allowbreak{}SIM\_\allowbreak{}APP\_\allowbreak{}v2/\allowbreak{}docs/\allowbreak{}v0.\allowbreak{}3.\allowbreak{}6-\allowbreak{}beta/\allowbreak{}THEORY\_\allowbreak{}COMPOSITION\_\allowbreak{}CONTRACT.\allowbreak{}md}.}
Monte Carlo / Sobol runs in the companion app explore this band; they do
not by themselves validate microscopic mechanisms.

\hypertarget{phonon-nuclear-coupling-strength}{%
\paragraph{2.8.1 Phonon-Nuclear Coupling
Strength}\label{phonon-nuclear-coupling-strength}}

A critical mechanism for LENR is the coupling between lattice phonons
and nuclear degrees of freedom. This coupling enables energy transfer
from the lattice to nuclear systems, potentially facilitating reactions
that would otherwise be energetically forbidden.

\textbf{Coupling Constant Definition:}

The phonon-nuclear coupling strength can be expressed as:

\[
g_{\text{ph-nuc}} = \frac{\partial V_{\text{nuclear}}}{\partial u} \cdot \sqrt{\frac{\hbar}{2M\omega_{\text{ph}}}}
\]

where: - \(V_{\text{nuclear}}\) is the nuclear potential - \(u\) is the
lattice displacement - \(M\) is the nuclear mass -
\(\omega_{\text{ph}}\) is the characteristic phonon frequency

\textbf{Calculation for Deuterons in Palladium:}

For deuterons in Pd lattice: - \(M = 3.344 \times 10^{-27}\) kg
(deuteron mass) - \(\omega_{\text{ph}} \approx 10^{13}\) rad/s (optical
phonon frequency in Pd-D) - \(\partial V / \partial u \approx 1-10\)
eV/Å (estimated from lattice strain effects)

\[
g_{\text{ph-nuc}} \approx (1-10) \times \sqrt{\frac{1.055 \times 10^{-34}}{2 \times 3.344 \times 10^{-27} \times 10^{13}}} \approx 0.1-1 \text{ eV/Å}
\]

\textbf{Energy Transfer Mechanism:}

The phonon-mediated energy transfer rate is given by:

\[
\Gamma_{\text{transfer}} = \frac{2\pi}{\hbar} |g_{\text{ph-nuc}}|^2 \rho(\omega) n_{\text{ph}}(\omega)
\]

where \(\rho(\omega)\) is the phonon density of states and
\(n_{\text{ph}}(\omega) = (e^{\hbar\omega/k_B T} - 1)^{-1}\) is the
Bose-Einstein distribution.

\textbf{Multi-Phonon Processes:}

For reactions requiring energy \(\Delta E \sim 1\) MeV, coherent
multi-phonon processes are necessary:

\[
P_{\text{multi-phonon}} \propto \left(\frac{g_{\text{ph-nuc}}}{\hbar \omega_{\text{ph}}}\right)^{N} \quad \text{where } N = \frac{\Delta E}{\hbar \omega_{\text{ph}}}
\]

For \(\Delta E \sim 1\) MeV and \(\hbar \omega_{\text{ph}} \sim 50\)
meV, \(N \sim 2 \times 10^7\) phonons must act coherently. This
highlights the importance of quantum coherence in LENR systems, as
discussed in Section 2.5.1.

\textbf{Critical Assessment:} The probability of
\(N \sim 2 \times 10^7\) phonons acting coherently through independent
coupling scales as \((g_{\text{ph-nuc}}/\hbar\omega_{\text{ph}})^N\),
which for any coupling ratio less than unity yields a probability
indistinguishable from zero (e.g.,
\(0.99^{2 \times 10^7} \sim 10^{-87,000}\)). This effectively rules out
the naive multi-phonon pathway for bridging the MeV energy gap.

However, several alternative mechanisms may circumvent this limitation:
1. \textbf{Stimulated phonon cascades:} Analogous to stimulated emission
in lasers, where each phonon absorption event increases the probability
of the next, creating a cooperative avalanche rather than independent
events. 2. \textbf{Coherent energy transfer via off-shell
intermediates:} Hagelstein's model proposes that energy transfer between
lattice and nuclear degrees of freedom proceeds through virtual
(off-shell) intermediate states, analogous to Raman scattering, where
the full \(\Delta E\) need not be assembled simultaneously. 3.
\textbf{Phonon-polariton condensation:} If the phonon system achieves a
coherent state (analogous to a phonon laser or Bose-Einstein condensate
of phonon-polaritons), the energy transfer becomes a collective quantum
process with fundamentally different statistics than independent phonon
coupling.

The resolution of the MeV energy gap through lattice-nuclear coupling
remains the central open theoretical challenge in LENR. See Section 6.1
for further discussion.

\textbf{Implications:}

\begin{enumerate}
\def\labelenumi{\arabic{enumi}.}
\tightlist
\item
  The coupling strength determines the rate of energy exchange between
  lattice and nuclear degrees of freedom
\item
  High phonon occupation numbers (achieved at elevated temperatures or
  through external driving) enhance energy transfer
\item
  Coherent phonon states can dramatically increase the effective
  coupling by enabling cooperative multi-phonon processes
\end{enumerate}

\begin{center}\rule{0.5\linewidth}{0.5pt}\end{center}

\hypertarget{computational-methods}{%
\subsection{3. Computational Methods}\label{computational-methods}}

\textbf{Methods epistemic map (peer-review).} Four layers are kept
distinct throughout this section and §5A: (i) \textbf{exploratory
models} --- labelled hypotheses (Hagelstein / TSC / Widom--Larsen /
trapped-state / three-body) with stated would-close conditions; (ii)
\textbf{conjectures / SPEC} --- mechanism forms not yet licensed by
governed data; (iii) \textbf{Lean theorems} --- identities, non-closure,
and BOUNDED / form- or class-DISPROVEN / proved-GAP results under
\emph{stated} caps or model classes (Lean 4 without mathlib; no
\texttt{sorry}; one labelled bridge axiom); (iv) \textbf{DEF
bookkeeping} --- clamps, bands, detector-flag mirrors, expected-yield
accounting that do not allocate MeV remainder, derive γ suppression,
close a rate budget, supply the EC threshold, or account for integrated
\(^4\mathrm{He}\). Simulation outputs remain claim-tiered sensitivity
analyses, never empirical posteriors. Affirmative closure of any
\texttt{EQ-OPEN-*} requires governed lab evidence.

\begin{quote}
\textbf{Implementation note (public beta, v0.3.7-beta):} Locked live
claim tiers --- \textbf{Beta:} WKB/Gamow tunneling
(\texttt{ENG-WKB-01}); \textbf{Approximate:} Rayleigh--Plesset /
Keller--Miksis bubble dynamics with KM auto-select when peak Mach
\(\ge 0.3\) (\texttt{ENG-RP-KM-01}). The live VPS runs the in-repo
\texttt{physics\_engine} pipeline via \texttt{/api/physics/*} and
\texttt{/api/unified/*}. Outputs carry provenance
(\texttt{physics\_engine} vs schematic fallback). \textbf{Aspirational
(NOT live):} COMSOL-class FEM, self-consistent Poisson--Schrödinger, and
coupled continuum solvers ({\ttfamily ENG-\allowbreak{}ASPIRATIONAL-\allowbreak{}FEM-\allowbreak{}01}) are
research methods only --- never public-beta numerics. Lean
\texttt{PoissonSchrodinger} / IR \texttt{EQ-PS-*} encode residual,
domain, and non-deployment honesty only; they do \textbf{not} prove LENR
or close any \texttt{EQ-OPEN-*}.
\end{quote}

\hypertarget{poissonschruxf6dinger-coupled-dynamics}{%
\subsubsection{3.1 Poisson--Schrödinger Coupled
Dynamics}\label{poissonschruxf6dinger-coupled-dynamics}}

\textbf{Honesty label:} Continuum FEM / self-consistent
Poisson--Schrödinger below is an \textbf{aspirational modeling roadmap}
({\ttfamily ENG-\allowbreak{}ASPIRATIONAL-\allowbreak{}FEM-\allowbreak{}01}). It is \textbf{not} executed on the
public beta, \textbf{not} a COMSOL production path in this release, and
\textbf{not} a discharge of any \texttt{EQ-OPEN-*}. Formalization locks
parent continuum DEF \texttt{EQ-PS-01}, discrete residual algebra,
permittivity domain, explicit tolerance axiom, and non-deployment flags
(\texttt{EQ-PS-POISSON-01}, \texttt{EQ-PS-SCHROD-01},
\texttt{EQ-PS-COUPLE-01}, \texttt{EQ-PS-DOM-01}, \texttt{EQ-PS-TOL-01},
{\ttfamily EQ-\allowbreak{}PS-\allowbreak{}NONDEPLOY-\allowbreak{}01}).

\begin{enumerate}
\def\labelenumi{\arabic{enumi}.}
\item
  \textbf{Poisson's Equation} (parent continuum DEF \texttt{EQ-PS-01}):
  \[
  \nabla^{2}\phi = -\frac{\rho}{\epsilon_0}
  \quad\text{or more generally}\quad
  \nabla\cdot\bigl(\varepsilon\nabla\phi\bigr)=-\rho,
  \] where \(\phi\) is the electric potential, \(\rho\) is the charge
  density, and \(\varepsilon\) is a (possibly spatially varying)
  permittivity. Domain: permittivity must be positive wherever the
  formula is claimed (\texttt{EQ-PS-DOM-01}).
\item
  \textbf{Schrödinger's Equation:} \[
  -\frac{\hbar^{2}}{2m}\nabla^{2}\psi + V(\mathbf{r})\psi = E\psi,
  \] which governs the quantum state \(\psi\) under potential
  \(V(\mathbf{r})\). Here \(E\) is the energy eigenvalue
  (\texttt{E\_eigen}), not the electric field (\texttt{E\_field}).
\item
  \textbf{Self-consistent coupling (roadmap):} iterate Poisson \(\to\)
  update \(V\leftarrow V_{\mathrm{ext}}-e\phi\) \(\to\) Schrödinger
  \(\to\) update \(\rho\) from \(|\psi|^{2}\) until residual
  \(r=\max(r_{\mathrm{P}},r_{\mathrm{S}})\) vanishes in the discrete
  bookkeeping sense (\texttt{EQ-PS-COUPLE-01}), or \(r\le\tau\) under an
  \textbf{explicit} tolerance axiom (\texttt{EQ-PS-TOL-01} --- never
  silently promoted to exact \(r=0\)). Continuum existence / FEM
  convergence are \textbf{not} theorems of this pilot.
\item
  \textbf{Simulation Details (research targets, not VPS execution)}

  \begin{itemize}
  \tightlist
  \item
    \textbf{Software target:} COMSOL Multiphysics and custom Python
    scripts (NumPy, SciPy) for finite-element and finite-difference
    methods --- aspirational.\\
  \item
    \textbf{Grid Size:} \(\sim 0.1\,\text{nm}\).\\
  \item
    \textbf{Time Steps:} \(\sim 1\,\text{fs}\) (time-dependent
    extensions).\\
  \item
    \textbf{Illustrative expectations:} Electric fields in interfacial
    regions can exceed \(10^{8}\,\text{V/m}\); wavefunction solutions
    can show elevated electron density near protons. Any link to
    neutron-formation pathways remains claim-tiered and speculative ---
    \textbf{not} nuclear-rate proof.
  \end{itemize}
\end{enumerate}

\begin{center}\rule{0.5\linewidth}{0.5pt}\end{center}

\hypertarget{bubble-collapse-dynamics}{%
\subsubsection{3.2 Bubble Collapse
Dynamics}\label{bubble-collapse-dynamics}}

\textbf{Honesty label:} Cavitation--collapse in this framework is
\textbf{modelled enhancement / contested hydro}. Rayleigh--Plesset /
Keller--Miksis trajectories, surface-proximity factors, and collapse UQ
percentiles are companion-software sensitivity explorations under
claim-tier governance. They are \textbf{not} laboratory proof of nuclear
rates, \textbf{not} a discharge of any \texttt{EQ-OPEN-*}, and
\textbf{not} a physical-necessity theorem. Lean
\texttt{LenrFormal.CollapseBounds} encodes only clamp / bound / UQ /
RP-domain / hard-core CR-cap bookkeeping (\texttt{EQ-BBL-LIM-01},
\texttt{EQ-BBL-BOUND-01}, \texttt{EQ-BBL-UQ-01},
\texttt{EQ-BBL-RP-DOM-01}, \texttt{EQ-BBL-CR-CAP-01})---including the
pointwise clamp \(\min(\mathrm{raw},\mathrm{bound})\)---not
cavitation-dynamics theorems that prove LENR. \texttt{EQ-BBL-NR-01}
proves {\ttfamily hydro\_\allowbreak{}enhancement\_\allowbreak{}does\_\allowbreak{}not\_\allowbreak{}entail\_\allowbreak{}nuclear\_\allowbreak{}rate}: a
rate scenario containing a hydro/collapse enhancement or scenario prior
is definitionally \textbf{not} a nuclear-rate proof. WH-ENP-01 /
OpenObligations prove \textbf{non-closure} implications under stated
deficits---not affirmative LENR
({\ttfamily docs/\allowbreak{}FORMALIZATION/\allowbreak{}FORMAL\_\allowbreak{}PROOF\_\allowbreak{}STATUS.\allowbreak{}md}).

\textbf{Context:} During electrolysis, microbubbles can form and then
collapse, concentrating energy {[}26{]}. Pyyhtiä et al.~(2026)
(\texttt{pyyhtia2026}) provide OA experimental support that HER-driven
nanobubble cavitation damages CR-39 and confounds SPAWAR-class
pit=nuclear detector interpretations (pits without Pd and/or D₂O under
suitable geometry; HIFU control; shallow polymer pits vs deep α
tracks)---\textbf{not} LENR-rate proof; \texttt{EQ-OPEN-RAD-01} stays
OPEN. Ledger IDs: {\ttfamily ENG-\allowbreak{}BUBBLE-\allowbreak{}FUSION-\allowbreak{}01}, \texttt{ENG-RP-KM-01}
(claim tier \textbf{approximate}; do not treat collapse numerics as
nuclear-rate proof). Collapse UQ in the companion engine
(\texttt{collapse\_uq.py}, Sobol) reports modelled percentiles; campaign
grids that stress cavitation/collapse cells remain model output only
(see §3.8).

\begin{enumerate}
\def\labelenumi{\arabic{enumi}.}
\item
  \textbf{Modified Rayleigh--Plesset Equation} (incompressible /
  low-Mach baseline): \[
  R\,\ddot{R} + \frac{3}{2}\dot{R}^{2} = \frac{1}{\rho}\left(P_{\text{bubble}} - P_{\infty} - \frac{2\sigma_{\text{surf}}}{R} - \frac{4\mu\,\dot{R}}{R} + \frac{\epsilon_0 E^{2}}{2}\right),
  \] where \(R\) is bubble radius, \(\rho\) fluid density,
  \(\sigma_{\text{surf}}\) surface tension, \(\mu\) viscosity, and \(E\)
  external electric field.

  \emph{Note:} The electric pressure term
  \(\epsilon_0 E^2 / 2 \approx 4.4 \times 10^4\) Pa at \(E = 10^8\) V/m
  is negligible compared to collapse pressures of \textasciitilde GPa.
  This term primarily influences bubble initiation and early-stage
  dynamics rather than peak collapse conditions. Rayleigh--Plesset is
  \textbf{domain-limited} when peak Mach numbers are not \(\ll 1\).
\item
  \textbf{Keller--Miksis compressibility correction (companion software,
  v0.3.7-beta):} When liquid compressibility / acoustic radiation
  matters, the Unified and Bubble Fusion paths may use
  {\ttfamily dynamics\_\allowbreak{}mode=\allowbreak{}auto} and select Keller--Miksis when peak Mach
  \(\ge 0.3\). Explicit overrides (\texttt{legacy} /
  \texttt{rayleigh\_plesset} / \texttt{keller\_miksis}) remain
  available. Making KM the unconditional default, real-gas EOS, and FEM
  remain deferred ({\ttfamily DEFERRED\_\allowbreak{}MULTIPHYSICS\_\allowbreak{}FIDELITY.\allowbreak{}md}). LaTeX
  stub: {\ttfamily docs/\allowbreak{}v0.\allowbreak{}3.\allowbreak{}7-\allowbreak{}beta/\allowbreak{}latex/\allowbreak{}eng\_\allowbreak{}rp\_\allowbreak{}km\_\allowbreak{}01.\allowbreak{}tex}.
\item
  \textbf{Provenance honesty:} UI and export envelopes must distinguish
  {\ttfamily provenance=\allowbreak{}physics\_\allowbreak{}engine} ODE trajectories from schematic /
  local fallbacks. Schematic output must never be presented as validated
  engine numerics.
\item
  \textbf{Simulation Outcomes:}

  \begin{itemize}
  \tightlist
  \item
    \textbf{Temperature Peaks:} Localized temperatures can exceed 10⁵ K
    in extreme regimes. \emph{Note:} Peak temperatures depend strongly
    on the collapse regime. Single-bubble sonoluminescence in water
    typically achieves \textasciitilde10⁴ K, while the 10⁵ K figure
    applies to extreme multi-bubble collapse or laser-driven cavitation
    conditions. The relevant regime for LENR electrolysis systems should
    be specified when applying these estimates.\\
  \item
    \textbf{Pressure Increases:} Collapsing bubbles can produce
    pressures of multiple GPa. These collapse conditions may alter
    interfacial transport and material microstructure, but do not by
    themselves establish an appreciable nuclear reaction rate.
  \end{itemize}
\end{enumerate}

\begin{center}\rule{0.5\linewidth}{0.5pt}\end{center}

\hypertarget{heterogeneous-nucleation-at-surfaces}{%
\subsubsection{3.3 Heterogeneous Nucleation at
Surfaces}\label{heterogeneous-nucleation-at-surfaces}}

Bubble formation in LENR systems occurs predominantly at lattice
surfaces (heterogeneous nucleation), where cavities, defects, and
wettability reduce the energy barrier compared to homogeneous nucleation
in the bulk.

\begin{enumerate}
\def\labelenumi{\arabic{enumi}.}
\tightlist
\item
  \textbf{Contact angle factor:} For a surface with contact angle θ, the
  heterogeneous nucleation barrier is reduced by f(θ) = (2 + cos θ)(1 −
  cos θ)²/4. For θ = 90°, f = 0.5; for complete wetting θ → 0°, f → 0.
\item
  \textbf{Critical radius (Young--Laplace):} r* = 2γ/ΔP, where γ is
  liquid--vapor surface tension and ΔP is the pressure difference
  between vapor and liquid. At surfaces, geometry and contact angle
  modify the effective critical radius.
\item
  \textbf{Gibbs free energy barrier:} ΔG*\_het = ΔG*\_hom × f(θ), with
  ΔG*\_hom = (16π γ³)/(3 ΔP²). Lower f(θ) at hydrophobic or cavity sites
  promotes nucleation.
\item
  \textbf{Waiting time and active site probability:} Nucleation rate
  depends on exp(−ΔG*/(kT)); statistical waiting times and site
  activation probabilities follow from classical nucleation theory and
  supersaturation.
\end{enumerate}

\begin{center}\rule{0.5\linewidth}{0.5pt}\end{center}

\hypertarget{bubblesurface-coupling}{%
\subsubsection{3.4 Bubble--Surface
Coupling}\label{bubblesurface-coupling}}

The physical coupling between a cavitation bubble and the lattice
surface determines equilibrium standoff, attachment/detachment, and
asymmetric collapse.

\begin{enumerate}
\def\labelenumi{\arabic{enumi}.}
\tightlist
\item
  \textbf{Equilibrium distance:} Balance of Casimir (vacuum energy), van
  der Waals, and thermal forces sets the equilibrium separation between
  bubble and surface; standoff ratio γ = h/R (gap/radius) affects
  collapse symmetry.
\item
  \textbf{Van der Waals (sphere--plate):} Derjaguin approximation gives
  the force between a bubble and a flat surface; Hamaker constants for
  water--metal systems determine attraction at nanoscale gaps.
\item
  \textbf{Detachment radius:} Fritz-style relations describe bubble
  departure size when buoyancy and dynamic forces overcome attachment
  (surface tension, contact line).
\item
  \textbf{Microjet formation:} Asymmetric collapse near a rigid boundary
  produces a high-speed jet toward the surface; jet velocity and impact
  pressure scale with standoff and drive localized energy deposition.
  Any link from microjets to nuclear rates remains \textbf{modelled
  enhancement / contested hydro}, not an empirical LENR claim.
\item
  \textbf{Surface proximity enhancement:} For γ \textless{} 2, collapse
  pressure and temperature are amplified relative to spherical collapse
  in the hydro models; integrated software may apply surface-proximity
  factors inside claim-tiered rate \textbf{scenarios}. Those factors do
  not close \texttt{EQ-OPEN-RATE-01} or raise empirical claim tiers.
\end{enumerate}

\begin{center}\rule{0.5\linewidth}{0.5pt}\end{center}

\hypertarget{thermal-and-diffusion-boundary-layers}{%
\subsubsection{3.5 Thermal and Diffusion Boundary
Layers}\label{thermal-and-diffusion-boundary-layers}}

Thin layers adjacent to the lattice control where nucleation is favored
and how heat and mass are transported.

\begin{enumerate}
\def\labelenumi{\arabic{enumi}.}
\tightlist
\item
  \textbf{Thermal boundary layer:} Thickness δ\_th from Nusselt-number
  correlations (e.g., Nu ∝ Re\^{}a Pr\^{}b); determines the superheat
  profile and the region where liquid is hot enough to nucleate bubbles.
\item
  \textbf{Diffusion boundary layer:} Thickness δ\_D scales with Lewis
  number (Le = α/D); concentration gradients of dissolved gas (e.g., D₂)
  set supersaturation and nucleation driving force.
\item
  \textbf{Nucleation zone thickness:} The overlap of
  superheated/supersaturated regions defines where heterogeneous
  nucleation is most likely; material properties (Pd, Ni, D₂O) and flow
  conditions set typical scales (µm to tens of µm).
\item
  \textbf{Temperature and concentration profiles:} Erf-based or similar
  profiles describe T(y) and C(y) from surface to bulk; used in
  integrated surface-reaction and nucleation models.
\end{enumerate}

\begin{center}\rule{0.5\linewidth}{0.5pt}\end{center}

\hypertarget{acoustic-and-ultrasonic-forcing}{%
\subsubsection{3.6 Acoustic and Ultrasonic
Forcing}\label{acoustic-and-ultrasonic-forcing}}

External or flow-driven acoustic fields alter bubble dynamics and
nucleation in electrolytic and sonochemical LENR configurations.

\begin{enumerate}
\def\labelenumi{\arabic{enumi}.}
\tightlist
\item
  \textbf{Minnaert frequency:} Natural bubble resonance f₀ ∝
  (1/R)√(3κP/ρ); driving near f₀ amplifies radial oscillations and can
  promote rectified diffusion and collapse severity.
\item
  \textbf{Primary Bjerknes force:} Acoustic radiation force on a bubble
  in a pressure gradient drives it toward pressure antinodes (or nodes,
  depending on compressibility), concentrating bubbles in specific
  regions.
\item
  \textbf{Secondary Bjerknes force:} Bubble--bubble interaction force
  (e.g., attraction for similar phase); affects clustering and
  coalescence.
\item
  \textbf{Acoustic streaming:} Rayleigh-type streaming near boundaries
  creates steady circulation; velocity scales with
  (p\_a²/(ρc²))(L/δ\_v), influencing boundary layer and nucleation zone.
\item
  \textbf{Standing wave positions:} Node/antinode locations set where
  pressure and velocity amplitudes are maximal; used to place electrodes
  or nucleation sites.
\item
  \textbf{Acoustic nucleation enhancement:} Pressure oscillations lower
  the effective nucleation barrier (e.g., reduction of ΔG* in
  exp(−ΔG*/(kT))), increasing nucleation rate.
\item
  \textbf{Rectified diffusion:} Asymmetric gas flux over the oscillation
  cycle causes net gas uptake by bubbles, leading to growth and eventual
  collapse; important for sustained cavitation.
\end{enumerate}

\begin{center}\rule{0.5\linewidth}{0.5pt}\end{center}

\hypertarget{non-equilibrium-thermodynamics}{%
\subsubsection{3.7 Non-Equilibrium
Thermodynamics}\label{non-equilibrium-thermodynamics}}

LENR systems operate far from thermodynamic equilibrium, requiring a
proper treatment of energy flow, entropy production, and dissipative
structures. This section provides the thermodynamic foundation for
understanding energy transport in active LENR systems.

\hypertarget{energy-continuity-equation}{%
\paragraph{3.7.1 Energy Continuity
Equation}\label{energy-continuity-equation}}

The fundamental energy balance in a LENR system is governed by the
continuity equation:

\[
\frac{\partial \rho_E}{\partial t} + \nabla \cdot \mathbf{J}_E = \sigma_{\text{source}}
\]

where: - \(\rho_E\) is the local energy density {[}J/m³{]} -
\(\mathbf{J}_E\) is the energy current density {[}W/m²{]} -
\(\sigma_{\text{source}}\) is the volumetric source term from nuclear
reactions {[}W/m³{]}

For LENR systems, the source term includes:

\[
\sigma_{\text{source}} = R_{\text{nuc}} \cdot Q_{\text{reaction}} - \sigma_{\text{rad}} - \sigma_{\text{chem}}
\]

where \(R_{\text{nuc}}\) is the nuclear reaction rate,
\(Q_{\text{reaction}}\) is the energy per reaction,
\(\sigma_{\text{rad}}\) accounts for radiative losses, and
\(\sigma_{\text{chem}}\) represents chemical energy contributions.

\hypertarget{entropy-production}{%
\paragraph{3.7.2 Entropy Production}\label{entropy-production}}

The local entropy production rate in LENR systems is:

\[
\dot{s} = \frac{1}{T}\left[\mathbf{J}_E \cdot \nabla\left(\frac{1}{T}\right) + \sum_i \mathbf{J}_i \cdot \nabla\left(\frac{\mu_i}{T}\right) + \sigma_{\text{source}}\left(1 - \frac{T_0}{T}\right)\right]
\]

where: - \(T\) is the local temperature - \(\mathbf{J}_i\) are species
fluxes - \(\mu_i\) are chemical potentials - \(T_0\) is a reference
temperature

The second law requires \(\dot{s} \geq 0\), constraining the allowed
energy flow patterns.

\hypertarget{heat-transport-equations}{%
\paragraph{3.7.3 Heat Transport
Equations}\label{heat-transport-equations}}

For practical LENR systems, heat transport follows:

\[
\rho c_p \frac{\partial T}{\partial t} = \nabla \cdot (\kappa \nabla T) + \sigma_{\text{source}}
\]

where: - \(\rho\) is the material density - \(c_p\) is the specific heat
capacity - \(\kappa\) is the thermal conductivity

\textbf{Boundary Conditions for Calorimetry:}

At the electrode-electrolyte interface: \[
-\kappa \left.\frac{\partial T}{\partial n}\right|_{\text{interface}} = h(T_{\text{surface}} - T_{\text{bulk}}) + \epsilon \sigma_{\text{SB}}(T_{\text{surface}}^4 - T_{\text{ambient}}^4)
\]

where \(h\) is the convective heat transfer coefficient, \(\epsilon\) is
emissivity, and \(\sigma_{\text{SB}}\) is the Stefan-Boltzmann constant.

\hypertarget{connection-to-experimental-measurements}{%
\paragraph{3.7.4 Connection to Experimental
Measurements}\label{connection-to-experimental-measurements}}

The excess power measured in calorimetry experiments is:

\[
P_{\text{excess}} = \int_V \sigma_{\text{source}} \, dV = P_{\text{out}} - P_{\text{in}} - P_{\text{chem}}
\]

Time-resolved measurements can be analyzed using:

\[
P_{\text{excess}}(t) = m c_p \frac{dT}{dt} + UA(T - T_{\text{ambient}}) - P_{\text{electrical}}
\]

where \(UA\) is the overall heat transfer coefficient-area product.

\textbf{Implications:}

\begin{enumerate}
\def\labelenumi{\arabic{enumi}.}
\tightlist
\item
  Non-equilibrium thermodynamics provides the framework for interpreting
  calorimetric data
\item
  Entropy production constraints help identify physically plausible
  energy flow patterns
\item
  Transient analysis can distinguish between nuclear and chemical heat
  sources based on temporal signatures
\end{enumerate}

\begin{center}\rule{0.5\linewidth}{0.5pt}\end{center}

\hypertarget{companion-software-simulation-evidence-v0.3.7-beta}{%
\subsubsection{3.8 Companion-software simulation evidence
(v0.3.7-beta)}\label{companion-software-simulation-evidence-v0.3.7-beta}}

Pinned, version-controlled model runs---not laboratory
experiments---anchor the quantitative claims in this revision.
Authoritative evidence lives under
{\ttfamily LENR\_\allowbreak{}SIM\_\allowbreak{}APP\_\allowbreak{}v2/\allowbreak{}datasets/\allowbreak{}reviewer\_\allowbreak{}packs/\allowbreak{}} and the paper
multi-format pack provenance file (\texttt{PROVENANCE.md}).

\begin{enumerate}
\def\labelenumi{\arabic{enumi}.}
\tightlist
\item
  \textbf{Reviewer packs:} \texttt{smoke/} (CI smoke) and
  \texttt{q\_campaign\_demo/} (larger campaign: 150 experiments / 132
  theory runs in the recorded manifest). Schema
  {\ttfamily lenr-\allowbreak{}reviewer-\allowbreak{}pack-\allowbreak{}v1}. Disclaimer on every pack: model output
  only; additional runs do \textbf{not} raise empirical claim tiers or
  close \texttt{EQ-OPEN-*}.
\item
  \textbf{Cavitation / collapse campaign pack:}
  {\ttfamily LENR\_\allowbreak{}SIM\_\allowbreak{}APP\_\allowbreak{}v2/\allowbreak{}datasets/\allowbreak{}reviewer\_\allowbreak{}packs/\allowbreak{}cavitation\_\allowbreak{}collapse\_\allowbreak{}campaign/\allowbreak{}}
  (Phase A; on \texttt{dev} via \#299). Model-only bubble-lattice
  host×loading×drive×dynamics\_mode cells with collapse UQ; complements
  engine modules (\texttt{collapse\_uq.py}, RP/KM paths). Contested
  Taleyarkhan sonofusion literature is an anchor, not laboratory proof.
  Pyyhtiä et al.~2026 (\texttt{pyyhtia2026}) sits alongside as an
  electrolytic CR-39 cavitation / detector-confounder peer---not a
  nuclear-rate proof.
\item
  \textbf{Detector / weak screen campaign pack:}
  {\ttfamily LENR\_\allowbreak{}SIM\_\allowbreak{}APP\_\allowbreak{}v2/\allowbreak{}datasets/\allowbreak{}reviewer\_\allowbreak{}packs/\allowbreak{}detector\_\allowbreak{}weak\_\allowbreak{}screen/\allowbreak{}}
  (WS7; on \texttt{dev} via \#331). Model-only cells for
  detector/weak/screening inventory; dual
  \texttt{score}/\texttt{final\_score} honesty preserved; does
  \textbf{not} close \texttt{EQ-OPEN-*} or set \texttt{lenr\_proved}.
\item
  \textbf{Campaign + UQ / avenue ranking:} Rate and collapse UQ feed
  decision-support avenue ranking at
  {\ttfamily LENR\_\allowbreak{}SIM\_\allowbreak{}APP\_\allowbreak{}v2/\allowbreak{}docs/\allowbreak{}v0.\allowbreak{}3.\allowbreak{}9-\allowbreak{}beta/\allowbreak{}AVENUE\_\allowbreak{}RANKING.\allowbreak{}md} (Phase
  B; on \texttt{dev} via \#301). Ranking prioritizes which lab avenues
  look informative next; it does \textbf{not} set \texttt{lenr\_proved}
  or close OPEN ids. Companion Lean \texttt{Decision} / IR
  \texttt{EQ-UQ-BF-01} ExpHom identities are \textbf{algebraic honesty
  floors only}---not empirical Bayes factors or laboratory posteriors.
\item
  \textbf{Theory weight surfaces:}
  {\ttfamily LENR\_\allowbreak{}SIM\_\allowbreak{}APP\_\allowbreak{}v2/\allowbreak{}docs/\allowbreak{}v0.\allowbreak{}3.\allowbreak{}9-\allowbreak{}beta/\allowbreak{}THEORY\_\allowbreak{}WEIGHT\_\allowbreak{}SURFACES.\allowbreak{}md}
  (WS7; on \texttt{dev} via \#331) inventories
  {\ttfamily CLAIM\_\allowbreak{}TIER\_\allowbreak{}MULTIPLIER}, oracle priors, and avenue blends as
  \textbf{Approximate / Aspirational} governance weights---not measured
  lab posteriors.
\item
  \textbf{Campaign invariants:}
  {\ttfamily formalization/\allowbreak{}predictions/\allowbreak{}campaign\_\allowbreak{}invariants.\allowbreak{}json} and
  pack-embedded invariants enforce
  {\ttfamily open\_\allowbreak{}questions.\allowbreak{}lenr\_\allowbreak{}proved\ =\allowbreak{}\ false}, the five OPEN IDs,
  two-body \(\log_{10}N_{\mathrm{events}}\approx-40\), and
  {\ttfamily contract\_\allowbreak{}version\ =\allowbreak{}\ v0.\allowbreak{}3.\allowbreak{}7-\allowbreak{}beta-\allowbreak{}science-\allowbreak{}gate-\allowbreak{}1}.
\item
  \textbf{Export envelope:} {\ttfamily /\allowbreak{}api/\allowbreak{}unified/\allowbreak{}export} uses
  {\ttfamily lenr-\allowbreak{}run-\allowbreak{}export-\allowbreak{}v1} with claim-tier / provenance fields
  required by the data contract.
\item
  \textbf{Paper benchmarks:}
  {\ttfamily backend/\allowbreak{}tests/\allowbreak{}test\_\allowbreak{}paper\_\allowbreak{}benchmarks/\allowbreak{}} locks selected theory
  numeric/regime targets in CI.
\item
  \textbf{Lean clamp / non-closure / bounds (not closure):}
  \texttt{CollapseBounds}
  ({\ttfamily EQ-\allowbreak{}BBL-\allowbreak{}LIM/\allowbreak{}BOUND/\allowbreak{}UQ/\allowbreak{}RP-\allowbreak{}DOM/\allowbreak{}CR-\allowbreak{}CAP-\allowbreak{}01}, \texttt{EQ-BBL-NR-01}
  hydro≠nuclear), \texttt{MeVRemainderCaps} (\texttt{EQ-WH-MEV-CAP-01};
  BOUNDED Q1), \texttt{RadiationBounds} (\texttt{EQ-WH-FC-02}
  class-DISPROVEN; \texttt{EQ-RAD-DET-01} DEF only), \texttt{RateBudget}
  (\texttt{EQ-RATE-2B-NC-01}, {\ttfamily EQ-\allowbreak{}RATE-\allowbreak{}CAP-\allowbreak{}NC-\allowbreak{}01/\allowbreak{}02},
  {\ttfamily EQ-\allowbreak{}RATE-\allowbreak{}CAP-\allowbreak{}GAP-\allowbreak{}01}), \texttt{WeakInteraction}
  (\texttt{EQ-EC-*}, {\ttfamily EQ-\allowbreak{}OPEN-\allowbreak{}EC-\allowbreak{}BOUND-\allowbreak{}01}; ≠ EC closure),
  \texttt{OpenQuestionBounds} / \texttt{YieldAccounting}
  ({\ttfamily EQ-\allowbreak{}OPEN-\allowbreak{}YIELD-\allowbreak{}NC-\allowbreak{}01}; instantaneous ≠ integrated),
  \texttt{Overlap} (\texttt{EQ-OVERLAP-01/02}; ensemble ≠ rate),
  \texttt{PoissonSchrodinger} (\texttt{EQ-PS-01}/\texttt{EQ-PS-TOL-01};
  aspirational FEM), and WH / OpenObligations non-closure theorems
  constrain bookkeeping; they do \textbf{not} discharge
  \texttt{EQ-OPEN-*}. Draft closure/falsification criteria (still OPEN)
  live at
  {\ttfamily LENR\_\allowbreak{}SIM\_\allowbreak{}APP\_\allowbreak{}v2/\allowbreak{}docs/\allowbreak{}v0.\allowbreak{}3.\allowbreak{}9-\allowbreak{}beta/\allowbreak{}EQ\_\allowbreak{}OPEN\_\allowbreak{}CLOSURE\_\allowbreak{}RUNWAY.\allowbreak{}md}
  (Phase F; on \texttt{dev} via \#297).
\item
  \textbf{Neutron-detector honesty (RAD/RATE):} Companion
  {\ttfamily neutron\_\allowbreak{}rate\_\allowbreak{}bound} carries honesty class
  {\ttfamily detector\_\allowbreak{}ul\_\allowbreak{}not\_\allowbreak{}proof}. Cd-ratio / moderator-null /
  soft-X∧¬hard-γ PREDs (PRED-08/09/10) and contested Parkhomov curated
  rows are falsifiers / literature anchors---they do \textbf{not} close
  \texttt{EQ-OPEN-RAD-01} or \texttt{EQ-OPEN-RATE-01}, and ν/anti-ν
  detector-gas SPEC rows do not discharge \texttt{EQ-OPEN-EC-01}.
\item
  \textbf{Non-claims:} Larger \(N\) in a campaign does not convert
  speculative mechanism labels into empirical validation. Maturity axes
  (\texttt{model\_fidelity}, \texttt{empirical\_support},
  \texttt{claim\_tier}) remain independent
  ({\ttfamily PHYSICS\_\allowbreak{}MATURITY\_\allowbreak{}LEDGER.\allowbreak{}md}).
\end{enumerate}

\begin{center}\rule{0.5\linewidth}{0.5pt}\end{center}

\hypertarget{experimental-context-and-reported-correlations}{%
\subsection{4. Experimental Context and Reported
Correlations}\label{experimental-context-and-reported-correlations}}

\hypertarget{electrolytic-reactors}{%
\subsubsection{4.1 Electrolytic Reactors}\label{electrolytic-reactors}}

\begin{enumerate}
\def\labelenumi{\arabic{enumi}.}
\item
  \textbf{Setup:}

  \begin{itemize}
  \tightlist
  \item
    Palladium electrodes (\(>99.99\%\) purity).\\
  \item
    Heavy water (\(\text{D}_2\text{O}\)) electrolytes with lithium
    salts.\\
  \item
    Current densities between \(0.1\text{–}1\,\text{A/cm}^2\).\\
  \end{itemize}
\item
  \textbf{Measurements:}

  \begin{itemize}
  \tightlist
  \item
    \textbf{Calorimetry:} Input power (voltage \(\times\) current)
    vs.~output power via isoperibolic or flow calorimeters.\\
  \item
    \textbf{Mass Spectrometry:} Detection of helium-4, tritium, or other
    isotopes in off-gas.\\
  \item
    \textbf{Neutron Counters:} Helium-3 or \(\text{BF}_3\) detectors
    above background sensitivity. Detector upper limits in the companion
    software are honesty-classed {\ttfamily detector\_\allowbreak{}ul\_\allowbreak{}not\_\allowbreak{}proof};
    Cd-ratio / moderator-null / soft-X∧¬hard-γ PREDs (PRED-08/09/10) are
    branching falsifiers, not closures of \texttt{EQ-OPEN-RAD-01} /
    \texttt{EQ-OPEN-RATE-01}.\\
  \end{itemize}
\item
  \textbf{Reported outcomes (disputed; claim-tier: empirical
  correlation):}

  \begin{itemize}
  \tightlist
  \item
    Excess heat relative to electrical input has been reported in some
    electrolytic campaigns (often cited near tens of percent);
    replication and calorimetry systematics remain contested.\\
  \item
    Helium-4 and tritium signals have been reported in a subset of
    studies; collection efficiency and blank controls vary.\\
  \item
    Low-level neutron signals are occasionally reported above
    background; they are not established as routine LENR signatures, and
    a Cd-wrapped count is \textbf{not} neutron proof.
  \end{itemize}
\end{enumerate}

\begin{center}\rule{0.5\linewidth}{0.5pt}\end{center}

\hypertarget{gaseous-reactors}{%
\subsubsection{4.2 Gaseous Reactors}\label{gaseous-reactors}}

\begin{enumerate}
\def\labelenumi{\arabic{enumi}.}
\item
  \textbf{Setup:}

  \begin{itemize}
  \tightlist
  \item
    Nickel powders with high surface area.\\
  \item
    Hydrogen gas at 1--10 atm.\\
  \item
    Operating temperatures of 300--800 K.\\
  \end{itemize}
\item
  \textbf{Measurements:}

  \begin{itemize}
  \tightlist
  \item
    \textbf{Calorimetry:} Monitoring of potential excess heat.\\
  \item
    \textbf{Spectroscopy:} Raman, infrared, or other techniques to
    detect Rydberg Matter or ultra-dense hydrogen states.\\
  \item
    \textbf{Elemental Analysis:} Post-run detection of transmutation
    products such as copper and zinc in nickel.\\
  \end{itemize}
\item
  \textbf{Reported outcomes (disputed; claim-tier: empirical
  correlation):}

  \begin{itemize}
  \tightlist
  \item
    \textbf{Anomalous Heat:} High-power excess claims (including
    kW-scale reports) exist in the literature but lack community
    consensus on calorimetry.\\
  \item
    \textbf{RM Detection:} Spectroscopic signatures have been
    interpreted as consistent with Rydberg Matter in some studies;
    alternative assignments exist.\\
  \item
    \textbf{Transmutations:} Post-run elemental shifts on nickel
    surfaces are reported in some campaigns; contamination and
    analytical artifacts must be excluded case-by-case.
  \end{itemize}
\end{enumerate}

\begin{center}\rule{0.5\linewidth}{0.5pt}\end{center}

\hypertarget{results-and-discussion}{%
\subsection{5. Results and Discussion}\label{results-and-discussion}}

\hypertarget{computationalexperimental-correlation}{%
\subsubsection{5.1 Computational--Experimental
Correlation}\label{computationalexperimental-correlation}}

\begin{itemize}
\tightlist
\item
  \textbf{Tunneling Enhancements:} Calculations show up to 10\^{}3 to
  10\^{}5-fold increases in tunneling probability under combined RM,
  halo electron, and high-field environments.\\
\item
  \textbf{Neutron Formation:} Scenario neutron-rate bands in the
  companion software are claim-tiered sensitivity explorations under
  {\ttfamily detector\_\allowbreak{}ul\_\allowbreak{}not\_\allowbreak{}proof} honesty; they may be \emph{compared}
  to occasionally reported low-level neutron signals but do \textbf{not}
  validate those signals, close \texttt{EQ-OPEN-RAD-01} /
  \texttt{EQ-OPEN-RATE-01}, or establish routine LENR neutron
  production.
\end{itemize}

\hypertarget{statistical-mechanics-of-reaction-rates}{%
\paragraph{5.1.1 Statistical Mechanics of Reaction
Rates}\label{statistical-mechanics-of-reaction-rates}}

A rigorous treatment of LENR reaction rates requires proper statistical
mechanics, accounting for thermal distributions, enhancement factors,
and threshold conditions.

\textbf{Reaction Rate Formula:}

The nuclear reaction rate per unit volume for two reacting species is:

\[
R = \frac{n_1 n_2}{1+\delta_{12}} \langle \sigma v \rangle_{\text{thermal}} \times f_{\text{enhancement}} \times P_{\text{threshold}}
\]

where \(\delta_{12} = 1\) for identical reacting species (e.g., D+D
self-fusion, giving the familiar \(\tfrac{1}{2}n^2\) form used by
{\ttfamily reaction\_\allowbreak{}rate\_\allowbreak{}engine.\allowbreak{}py:\allowbreak{}:\allowbreak{}compute\_\allowbreak{}two\_\allowbreak{}body\_\allowbreak{}rate}) and
\(\delta_{12} = 0\) for distinct species. This is the \emph{thermal
self-fusion} rate and is distinct from the Oppenheimer--Phillips capture
rate \(R_{\text{rate}} = \sigma_{\text{OP}}\Phi_n N\) (§2.7, a
flux-driven capture process) and the generic placeholder
\(R_{\text{rate}} = n\sigma_{\text{nuc}}v\) (§7); see the Glossary (§8)
note on \(R\) for how these three regimes relate.

where: - \(n_1, n_2\) are the number densities of reacting species
{[}m⁻³{]} - \(\langle \sigma v \rangle_{\text{thermal}}\) is the
thermally-averaged cross-section-velocity product {[}m³/s{]} -
\(f_{\text{enhancement}}\) is the cumulative enhancement factor from
screening, coherence, etc. - \(P_{\text{threshold}}\) is the probability
of exceeding threshold conditions

\textbf{Thermal Averaging:}

The thermal average is computed over the Maxwell-Boltzmann distribution:

\[
\langle \sigma v \rangle = \sqrt{\frac{8}{\pi \mu (k_B T)^3}} \int_0^\infty \sigma(E) \, E \, e^{-E/k_B T} \, dE
\]

where \(\mu\) is the reduced mass and \(\sigma(E)\) is the
energy-dependent cross-section.

For tunneling-dominated reactions:

\[
\sigma(E) \propto \frac{1}{E} \exp\left(-\sqrt{\frac{E_G}{E}}\right)
\]

where \(E_G = 2\mu c^2 (\pi \alpha Z_1 Z_2)^2\) is the Gamow energy (in
SI,
\(E_G = 2\mu \left[\pi Z_1 Z_2 e^2/(4\pi\varepsilon_0 \hbar)\right]^2\);
for D+D, \(E_G \approx 986\ \text{keV}\)).

\textbf{Enhancement Factor Decomposition:}

The total enhancement factor combines multiple mechanisms:

\[
f_{\text{enhancement}} = f_{\text{screening}} \times f_{\text{field}} \times f_{\text{coherence}} \times f_{\text{phonon}}
\]

Each factor has associated uncertainties that propagate through the
calculation.

\textbf{Error Propagation:}

The uncertainty in the reaction rate is:

\[
\frac{\delta R}{R} = \sqrt{\left(\frac{\delta n_1}{n_1}\right)^2 + \left(\frac{\delta n_2}{n_2}\right)^2 + \left(\frac{\delta \langle \sigma v \rangle}{\langle \sigma v \rangle}\right)^2 + \left(\frac{\delta f}{f}\right)^2}
\]

\textbf{Statistical Significance Criteria:}

For experimental validation, statistical significance is assessed using
Poisson statistics:

\[
S = \frac{n_{\text{observed}} - n_{\text{background}}}{\sqrt{n_{\text{background}}}}
\]

A significance of \(S \geq 5\) (5σ) is typically required for claiming
positive detection. For excess heat measurements:

\[
S_{\text{heat}} = \frac{Q_{\text{excess}}}{\sigma_Q}
\]

where \(\sigma_Q\) includes both statistical and systematic
uncertainties.

\textbf{Monte Carlo Validation:}

Reaction rate predictions are validated using Monte Carlo sampling: 1.
Sample input parameters from their uncertainty distributions 2.
Propagate through the full calculation chain 3. Build probability
distributions for output observables 4. Compare with experimental data
using likelihood ratios

\emph{Honesty on sample size (companion app, 2026-08-26):} default Sobol
UQ remains \textbf{512 samples, opt-in} (API cap 512). Campaign presets
that a reviewer can actually run are \texttt{literature\_grid}
(\textasciitilde16--32 designed cells), \texttt{all\_claims}
(\textasciitilde6 unified + one run per registered theory), and optional
\texttt{diverse} (\textasciitilde237). That corpus is on the order of
\textbf{hundreds of model-output rows}, not \(10^4\)--\(10^6\). The
larger Monte Carlo aspiration remains a sensitivity method, not a
completed campaign. See
{\ttfamily LENR\_\allowbreak{}SIM\_\allowbreak{}APP\_\allowbreak{}v2/\allowbreak{}docs/\allowbreak{}v0.\allowbreak{}3.\allowbreak{}6-\allowbreak{}beta/\allowbreak{}RESEARCHER\_\allowbreak{}REVIEW\_\allowbreak{}PACK.\allowbreak{}md}.
Extra runs do not close §5A or prove LENR.

\textbf{Implications:}

\begin{enumerate}
\def\labelenumi{\arabic{enumi}.}
\tightlist
\item
  Proper statistical treatment enables rigorous comparison between
  theory and experiment
\item
  Error propagation identifies which parameters dominate the uncertainty
  budget
\item
  Statistical significance criteria provide objective standards for
  experimental claims
\item
  Monte Carlo methods enable exploration of parameter space and
  identification of optimal conditions
\end{enumerate}

\hypertarget{reported-nuclear-product-correlations-contested}{%
\subsubsection{5.2 Reported Nuclear-Product Correlations
(contested)}\label{reported-nuclear-product-correlations-contested}}

\emph{Claim-tier note.} This subsection inventories \textbf{reported}
isotopic / heat correlations from the contested LENR corpus. It does
\textbf{not} validate nuclear product detection, close
\texttt{EQ-OPEN-YIELD-01} / \texttt{EQ-OPEN-RAD-01}, or flip
\texttt{lenr\_proved}.

\begin{itemize}
\tightlist
\item
  \textbf{Isotopic Shifts:} Reports of helium-4, tritium, and heavier
  isotopes are sometimes \emph{interpreted as} multi-step nuclear
  processes (e.g., Oppenheimer-Phillips); alternative chemical /
  contamination / detector-confounder explanations remain open, and
  Pyyhtiä 2026 (\texttt{pyyhtia2026}) strengthens the CR-39 cavitation
  confounder beside SPAWAR-class pit claims.\\
\item
  \textbf{Excess Heat Consistency:} Reported experimental outputs are
  consistent in magnitude with the hypothesized nuclear energy
  contributions once purely chemical avenues are discounted; this
  remains a disputed empirical correlation (claim tier:
  {\ttfamily empirical\_\allowbreak{}lenr\_\allowbreak{}correlation}), not an established result.
\end{itemize}

\hypertarget{detailed-energy-accounting}{%
\paragraph{5.2.1 Detailed Energy
Accounting}\label{detailed-energy-accounting}}

For D+D → ⁴He reactions, the total energy release must be properly
accounted for:

\[
E_{\text{released}} = E_{\text{binding}}(^4\text{He}) - 2E_{\text{binding}}(D) = 23.8 \text{ MeV}
\]

This energy must be distributed among multiple channels:

\begin{enumerate}
\def\labelenumi{\arabic{enumi}.}
\tightlist
\item
  \textbf{Lattice Phonons:}
  \(E_{\text{phonon}} = \sum_i \hbar \omega_i n_i\), where \(\omega_i\)
  are phonon frequencies and \(n_i\) are occupation numbers
\item
  \textbf{Electronic Excitations:}
  \(E_{\text{electronic}} = \sum_k E_k f_k\), where \(E_k\) are
  electronic energy levels and \(f_k\) are occupation probabilities
\item
  \textbf{Nuclear Recoil:} \(E_{\text{recoil}} = \frac{p^2}{2M}\), where
  \(p\) is the recoil momentum and \(M\) is the nuclear mass
\item
  \textbf{Radiation (if not suppressed):}
  \(E_\gamma = \hbar \omega_\gamma\), where \(\omega_\gamma\) is the
  photon frequency
\end{enumerate}

Energy conservation requires:

\[
E_{\text{nuclear}} = E_{\text{phonon}} + E_{\text{electronic}} + E_{\text{recoil}} + E_\gamma
\]

For LENR, the suppression of high-energy gamma emission suggests:

\[
E_\gamma \ll E_{\text{nuclear}}
\]

implying that, \emph{if} a lattice-phonon channel dominated, most energy
would be channeled into lattice and electronic degrees of freedom. That
\textbf{hypothesis} is offered as one claim-tiered scenario for reported
excess heat without commensurate high-energy radiation; it does
\textbf{not} establish that observed heat is nuclear, does \textbf{not}
close \texttt{EQ-OPEN-RAD-01}, and is not a validated LENR signature.

\hypertarget{energy-product-yield-reconciliation}{%
\paragraph{5.2.2 Energy-Product Yield
Reconciliation}\label{energy-product-yield-reconciliation}}

A self-consistency check between excess heat and nuclear product yields
is essential. For the primary channel D+D → ⁴He + 23.8 MeV:

\begin{itemize}
\tightlist
\item
  Energy per reaction:
  \(Q = 23.8 \text{ MeV} = 3.81 \times 10^{-12} \text{ J}\)
\item
  Reactions required for 1 W·hr:
  \(N = 3600 \text{ J} / 3.81 \times 10^{-12} \text{ J} = 9.4 \times 10^{14} \text{ reactions}\)
\item
  Expected ⁴He yield: \(9.4 \times 10^{14} \text{ atoms per W·hr}\)
\end{itemize}

\textbf{Comparison with experimental claims:}

Reported ⁴He yields in LENR experiments range from
\(10^{11}\)--\(10^{13}\) atoms/W·hr (Miles 1993, McKubre 2006). This is
a factor of 100--10,000 below the theoretical yield if all excess heat
originated from D+D → ⁴He. Several explanations are possible:

\begin{enumerate}
\def\labelenumi{\arabic{enumi}.}
\tightlist
\item
  \textbf{Detection efficiency:} ⁴He measurement via mass spectrometry
  has limited collection efficiency, particularly when helium diffuses
  into the lattice or escapes through the electrolyte.
\item
  \textbf{Multiple reaction channels:} Heat may also arise from channels
  with different products: D+D → T + p (Q = 4.03 MeV), D+D → ³He + n (Q
  = 3.27 MeV), or lattice-mediated transmutations that produce heat
  without helium.
\item
  \textbf{Incomplete energy release:} If some reaction energy goes into
  lattice defect creation or radiation that escapes calorimetric
  detection, the effective heat per ⁴He atom would be lower.
\end{enumerate}

The Miles heat-helium correlation---\((2.6 \pm 0.6) \times 10^{11}\) ⁴He
atoms per second per watt of excess heat---is within a factor of
\textasciitilde2 of the 23.8 MeV/⁴He prediction (\(2.63 \times 10^{11}\)
atoms/s/W). If the underlying measurements are accepted, this
correlation is \emph{consistent with} a dominant D+D → ⁴He channel; it
remains a contested empirical correlation (helium retention and
measurement caveats apply) rather than a validated mechanism. The
discrepancy in integrated yields (per W·hr) likely reflects time-varying
reaction rates and measurement integration periods rather than a
fundamental inconsistency.

\textbf{Tritium and neutron accounting:}

\begin{itemize}
\tightlist
\item
  \textbf{Tritium:} Reported yields of \(10^5\)--\(10^7\) atoms/hr,
  \emph{if taken at face value}, imply a branching ratio of
  \(\lesssim 10^{-9}\) relative to a putative ⁴He channel. That is a
  contested corpus pattern---not proof that D+D → T+p is suppressed by a
  lattice-phonon pathway.
\item
  \textbf{Neutrons:} Reported neutron fluxes are typically \(< 10^3\)/s,
  corresponding to a branching ratio of \(\lesssim 10^{-11}\) under the
  same conditional framing. Companion-software detector upper limits and
  PRED-08/09/10 falsifiers are bookkeeping
  ({\ttfamily detector\_\allowbreak{}ul\_\allowbreak{}not\_\allowbreak{}proof}); they do \textbf{not} define a
  LENR signature and do \textbf{not} close \texttt{EQ-OPEN-RAD-01} /
  \texttt{EQ-OPEN-RATE-01}. Radiation-suppression hypotheses remain OPEN
  (§5A.2 / §6.1).
\end{itemize}

\hypertarget{implications-for-lenr}{%
\subsubsection{5.3 Implications for LENR}\label{implications-for-lenr}}

\begin{itemize}
\tightlist
\item
  \textbf{Unified Model:} By merging condensed matter, quantum
  mechanics, and nuclear physics, the framework offers a self-consistent
  candidate interpretation of multiple reported LENR observations; it
  does not by itself establish that those observations are nuclear in
  origin.\\
\item
  \textbf{Improved Reproducibility:} Standardized protocols with
  precisely characterized materials and well-monitored conditions are
  starting to yield more consistent data.
\end{itemize}

\begin{center}\rule{0.5\linewidth}{0.5pt}\end{center}

\hypertarget{a.-open-theoretical-challenges-and-fundamental-questions}{%
\subsection{5A. Open Theoretical Challenges and Fundamental
Questions}\label{a.-open-theoretical-challenges-and-fundamental-questions}}

Occurrence under Miles-class heat--helium correlations is empirically
anchored in parts of the corpus; the open theoretical work is
\textbf{how} MeV energy partitions (Q1), \textbf{how} prompt γ is
suppressed (Q2), \textbf{when} the rate budget closes under a named
mechanism path (Q3), \textbf{how} an electron-capture energy source
might reach threshold (Q4), and \textbf{when} integrated
heat--\(^4\mathrm{He}\) yield closes (Q5). \textbf{All five obligations
remain OPEN.} Formal Lean postures (BOUNDED / form- or class-DISPROVEN /
proved GAP) under stated assumptions narrow mechanism \emph{forms}; they
do \textbf{not} close any row and do \textbf{not} set
\texttt{lenr\_proved}. This section formulates each challenge with
testable falsifiers and closure predicates. The companion app's
\textbf{open-questions virtual lab} (\texttt{open\_questions} on unified
/ bubble-lattice / theory-pipeline responses) encodes occurrence
(\texttt{occurrence.anchored}), when ({\ttfamily when.\allowbreak{}threshold\_\allowbreak{}crossed},
duty cycle, stimulation), and how ({\ttfamily how.\allowbreak{}mechanism\_\allowbreak{}paths}, FGR
parameters) under Oracle A (Lean identities and conditionals) plus
Oracle B (mechanism-path odds). Oracle B Bayes-factor algebra
(\texttt{EQ-UQ-BF-01} ExpHom cocycle in \texttt{LenrFormal.Decision}) is
an \textbf{algebraic honesty floor} under stated inputs---not an
empirical posterior---and theory-weight surfaces remain
Approximate/Aspirational inventory
({\ttfamily THEORY\_\allowbreak{}WEIGHT\_\allowbreak{}SURFACES.\allowbreak{}md}). It never prints
\texttt{LENR\_PROVED}. See
{\ttfamily docs/\allowbreak{}ARCHITECTURE/\allowbreak{}OPEN\_\allowbreak{}QUESTIONS\_\allowbreak{}VIRTUAL\_\allowbreak{}LAB\_\allowbreak{}2026-\allowbreak{}08-\allowbreak{}22.\allowbreak{}md},
{\ttfamily docs/\allowbreak{}ARCHITECTURE/\allowbreak{}MANY\_\allowbreak{}BODY\_\allowbreak{}FGR\_\allowbreak{}MODEL.\allowbreak{}md}, and
{\ttfamily docs/\allowbreak{}REVIEW\_\allowbreak{}2026/\allowbreak{}OPEN\_\allowbreak{}QUESTION\_\allowbreak{}CORPUS\_\allowbreak{}GAPS.\allowbreak{}md}.

\textbf{Working hypothesis WH-ENP-01} (claim tier
{\ttfamily speculative\_\allowbreak{}exploratory}; Lean
{\ttfamily LenrFormal.\allowbreak{}WorkingHypothesis}). Electron--neutron--proton
coupling in a deuteron cluster is treated as one object: a Hagelstein
coherent virtual-phonon cascade on Takahashi TSC 4D geometry, with a
derived harmonic \(\Gamma_{\mathrm{ph}}/\Gamma_{\gamma}\) bound and the
same-cluster FGR versus \(R_{\mathrm{obs}}=P_{\mathrm{W}}/(QV)\).
Independent multi-phonon \(g^N\) remains ruled out
(\texttt{EQ-NAIVE-GN-01}). Bulk Heisenberg--Euler vacuum polarisation is
a negative control (\texttt{EQ-QED-HE-NEG-01}: lab-field factor
\(\equiv 1\)). Dicke \(N^2\) is \textbf{not} the engine rate law
(\texttt{EQ-COLL-LIN-01}); the paper's \(\sim 6.6\,\mathrm{nm}\) domain
is larger than the phonon cooperation radius \(\lambda/2\pi\), so
\(N^2\) is not licensed (\texttt{EQ-WH-DICKE-01}). Two on-shell alphas
remove a \emph{primary} 24 MeV \(\gamma\) kinematically
(\texttt{EQ-TSC-CH-01}) but do not close the MeV-gap
(\texttt{EQ-TSC-STOP-01}, \texttt{EQ-WH-STOP-01}). Lean proves
conservation and \textbf{non-closure} (\texttt{EQ-WH-PART-01},
\texttt{EQ-WH-BOUND-01}, \texttt{EQ-WH-NONCLOSE-01}).
\texttt{EQ-OPEN-MEV-01}, \texttt{EQ-OPEN-RAD-01}, and
\texttt{EQ-OPEN-RATE-01} stay OPEN.

\textbf{WH-ENP-01 v2 derived bounds} (same claim tier; Lean
{\ttfamily LenrFormal.\allowbreak{}WorkingHypothesisBounds}; runtime blocks
\texttt{q2.franck\_condon}, \texttt{q1.coupling\_scan},
{\ttfamily dicke.\allowbreak{}n\_\allowbreak{}within\_\allowbreak{}radius}, \texttt{q3.closure\_demand}). Four
quantitative bounds sharpen the remainder without discharging it. (i) A
displaced-oscillator (Huang--Rhys, linear-coupling) Franck--Condon
tail---a labelled model assumption---at phonon order
\(N=Q/hf\approx 5.75\times10^{8}\) (\(f=10^{13}\,\mathrm{Hz}\),
\(hf\approx 41.4\,\mathrm{meV}\)) gives
\(\log_{10}P_N\le-7\times10^{5}\ll-11\) for every Huang--Rhys factor
\(S\le10^{4}\); the smallest \(S\) that reaches the \(10^{-11}\) target
is \(S_{\min}\approx0.9998\,N\), i.e.~a lattice relaxation energy
\(S_{\min}hf\approx23.79\,\mathrm{MeV}\approx Q\)
(\texttt{EQ-WH-FC-01}). (ii) The deuteron count inside the phonon
cooperation sphere is \(n_D V_{\mathrm{dicke}}\approx0.033<1\)
(\(n_D=6.8\times10^{28}\,\mathrm{m^{-3}}\), \(v_s=3070\,\mathrm{m/s}\);
model assumptions), so \(N_{\mathrm{coh}}^{\max}=0\) and Dicke \(N^{2}\)
is not licensed for even one coherent pair (\texttt{EQ-WH-DICKE-N-01}).
(iii) A conservation-constrained coupling scan with first-order fraction
\(\min(1,n_{\mathrm{coh}}g^{2})\) (model assumption) leaves the cited
row (\(g=10^{-4}\), \(n_{\mathrm{coh}}=0\)) at a remainder of
\(23.8\,\mathrm{MeV}\); rows with \(n_{\mathrm{coh}}>0\) are unlicensed
scenarios (\texttt{EQ-WH-PART-02}). (iv) For \(1\,\mathrm{W}\) in a
nominal \(1\,\mathrm{cm^{3}}\),
\(R_{\mathrm{obs}}\approx2.6\times10^{17}\,\mathrm{m^{-3}s^{-1}}\)
demands
\(P_{\mathrm{occ}}f_{\mathrm{nuc}}C_{\mathrm{lin}}/\tau\approx5.8\times10^{26}\,\mathrm{s^{-1}}\)
at the assumed 25 pm trap, above the engine cap product
\(1\cdot1\cdot10^{6}/10^{-20}\,\mathrm{s}=10^{26}\,\mathrm{s^{-1}}\)
(\texttt{EQ-WH-CAP-01}). Lean composes these into
\(\neg Q_1\wedge\neg Q_2\wedge\neg Q_3\) (\texttt{EQ-WH-NONCLOSE-02});
\texttt{lenr\_proved} stays false.

\hypertarget{a.1-the-mev-energy-gap-problem-q1}{%
\subsubsection{5A.1 The MeV Energy Gap Problem
(Q1)}\label{a.1-the-mev-energy-gap-problem-q1}}

The central theoretical challenge in LENR: how does energy from
meV-scale lattice phonons (\(\hbar\omega_{\text{ph}} \sim 50\) meV)
bridge to MeV-scale nuclear processes (\(\Delta E \sim 1\text{-}24\)
MeV)?

\textbf{Quantitative assessment of proposed pathways:}

\begin{enumerate}
\def\labelenumi{\arabic{enumi}.}
\item
  \textbf{Direct multi-phonon coupling (ruled out):} As shown in Section
  2.8.1, the probability of \(N \sim 2 \times 10^7\) independently
  coupled phonons acting coherently is \(\sim 10^{-87,000}\). This
  pathway is definitively excluded for independent coupling.
\item
  \textbf{Hagelstein coherent energy transfer:} Hagelstein and Chaudhary
  (2015) propose that lattice-nuclear energy exchange proceeds through
  virtual (off-shell) intermediate states, analogous to stimulated Raman
  scattering at nuclear energy scales. In this model, the full
  \(\Delta E\) need not be assembled simultaneously; instead, energy is
  transferred coherently through a chain of virtual phonon-nuclear
  transitions. The predicted coupling matrix element scales as
  \(|V_{\text{eff}}|^2 \propto n_{\text{ph}}^N / N!\) for stimulated
  processes, which is dramatically larger than the independent coupling
  estimate. \emph{Status:} A complete QFT derivation remains
  outstanding, but the model is internally consistent and makes testable
  predictions about the relationship between excess heat and phonon
  spectrum.
\item
  \textbf{Superradiant phonon cascades:} By analogy with Dicke
  superradiance in quantum optics, if \(N_{\text{atoms}}\) deuterons
  couple coherently to a phonon field, the collective emission rate
  scales as \(N_{\text{atoms}}^2\) rather than \(N_{\text{atoms}}\). For
  \(N \sim 10^4\) atoms in a coherent domain (\textasciitilde6.6 nm),
  this gives an enhancement of \(\sim 10^8\) over incoherent processes.
  \emph{Status:} Theoretically promising but requires explicit
  calculation of the collective coupling strength for nuclear
  transitions.
\item
  \textbf{Non-equilibrium BEC pathway:} As proposed in Section 2.6, if
  the deuteron subsystem reaches effective BEC conditions through phonon
  standing wave cooling, the entire condensate acts as a single quantum
  entity. Nuclear transitions of the condensate involve collective
  wavefunction overlap, potentially circumventing the single-pair
  tunneling bottleneck entirely. \emph{Status:} The quantitative
  estimate in Section 2.6 shows BEC is formally possible
  (\(T_{\text{eff}} < T_c^{\text{eff}}\)), but the nuclear transition
  rate of the condensate has not been calculated.
\end{enumerate}

\textbf{Conclusion on Q1 (\texttt{EQ-OPEN-MEV-01} stays OPEN; Lean
posture BOUNDED):} The MeV energy gap is not bridged. Under the engine's
stated channel caps (phonon, γ, recoil), Lean \texttt{MeVRemainderCaps}
/ \texttt{EQ-WH-MEV-CAP-01} proves the remainder stays strictly positive
(\(R_{\min}=23776\) keV at \(\mathrm{BR}_\gamma=0\); \(R_{\min}>0\) for
any stated \(\mathrm{BR}_\gamma\le0.998\)); the only capped route to
zero remainder requires \(\mathrm{BR}_\gamma\ge0.999\), which
contradicts radiation suppression
({\ttfamily q1\_\allowbreak{}closure\_\allowbreak{}under\_\allowbreak{}engine\_\allowbreak{}caps\_\allowbreak{}forces\_\allowbreak{}gamma\_\allowbreak{}route} +
{\ttfamily gamma\_\allowbreak{}route\_\allowbreak{}closure\_\allowbreak{}fails\_\allowbreak{}q2}). Collapse/ENC hydro channels
allocate 0 MeV remainder ({\ttfamily hydro\_\allowbreak{}claims\_\allowbreak{}keep\_\allowbreak{}q1\_\allowbreak{}open}). This
is a \textbf{conditional bound}, not closure. WH-ENP-01 ranks
Hagelstein+TSC above ruled-out \(g^N\) and above bulk Schwinger/HE, but
the derived energy remainder remains \(23.8\,\mathrm{MeV}\) until a
constrained phonon-nuclear matrix element is supplied
(\texttt{EQ-WH-PART-01}). TSC two-alpha kinematics and an
order-of-magnitude Pd CSDA range for \(23.8\,\mathrm{MeV}\) \(\alpha\)
are bookkeeping, not closure. Dicke \(N^2\) is not licensed at the paper
domain, and the counting bound is stricter: fewer than one deuteron
(\(\approx0.033\), model-assumption density and sound speed) sits inside
the \(\lambda/2\pi\approx49\,\mathrm{pm}\) cooperation sphere, so no
coherent pair is licensed (\texttt{EQ-WH-DICKE-N-01}). The quantitative
form of the gap for linear phonon coupling is the Franck--Condon
requirement \(S_{\min}\approx N\approx5.75\times10^{8}\): the lattice
relaxation energy would have to equal essentially the full \(Q\)
(\texttt{EQ-WH-FC-01}, displaced-oscillator model assumption). The
conservation-constrained coupling scan keeps the cited remainder at
\(23.8\,\mathrm{MeV}\); only the unlicensed corner
\(g=0.1,\ n_{\mathrm{coh}}=100\) would zero it (\texttt{EQ-WH-PART-02}).
\textbf{Falsifier / closure:} a named phonon→nuclear (or ENC) mechanism
with a constrained nuclear matrix element and occupation probability
that drives \(E_{\text{open remainder}}\) to 0, \textbf{plus}
experimental licensing. The critical next step remains an explicit,
experimentally constrained coupling matrix element.

\hypertarget{a.2-radiation-suppression-mechanism-q2}{%
\subsubsection{5A.2 Radiation Suppression Mechanism
(Q2)}\label{a.2-radiation-suppression-mechanism-q2}}

If D+D → ⁴He releases 23.8 MeV, standard nuclear physics predicts a 23.8
MeV gamma ray. LENR experiments consistently show excess heat with
minimal gamma radiation (upper bounds typically \(< 10^{-11}\) of what
conventional D+D fusion would produce).

\textbf{Analysis of proposed suppression mechanisms:}

\begin{enumerate}
\def\labelenumi{\arabic{enumi}.}
\item
  \textbf{Internal conversion to phonons:} The 23.8 MeV must fragment
  into \(\sim 5 \times 10^8\) phonons (at 50 meV each). A single nuclear
  transition fragmenting into \(10^8\) final-state quanta faces an
  enormous phase-space suppression. However, if the nuclear transition
  occurs within a coherent lattice domain, the final state is a
  collective lattice excitation rather than individual phonons. The
  branching ratio depends critically on the density of states at the
  nuclear transition energy, which is dramatically enhanced in a
  crystalline environment due to Umklapp processes and defect-mediated
  momentum transfer.
\item
  \textbf{Lattice-assisted momentum conservation:} In free space, D+D →
  ⁴He + γ is the only two-body channel that conserves both energy and
  momentum. In a crystal, the lattice can absorb arbitrary momentum (in
  units of reciprocal lattice vectors), opening the D+D → ⁴He + phonons
  channel. The branching ratio scales with the squared matrix element
  for lattice coupling relative to radiative coupling:
  \(\Gamma_{\text{phonon}}/\Gamma_{\gamma} \propto (g_{\text{lattice}}/\alpha)^2 \times \rho_{\text{lattice}}(E) / \rho_{\gamma}(E)\),
  where \(\rho\) are the respective densities of final states.
\item
  \textbf{Quantitative constraint:} For the observed
  \(\Gamma_{\text{phonon}}/\Gamma_{\gamma} > 10^{11}\), the lattice
  coupling must be at least
  \(\sim 10^{5.5} \times \alpha \approx 10^{3.2}\) times the
  electromagnetic coupling, or the lattice density of states must exceed
  the photon density of states by a comparable factor, or some
  combination thereof. Whether these conditions are achievable in real
  PdD lattices remains an open calculation.
\end{enumerate}

\textbf{Conclusion on Q2 (\texttt{EQ-OPEN-RAD-01} stays OPEN; Lean
posture BOUNDED / class-DISPROVEN):} A harmonic first-principles bound
gives \(\Gamma_{\mathrm{ph}}/\Gamma_{\gamma}=0\) because the Pd Debye
cutoff (\(\sim 24\,\mathrm{meV}\)) lies far below \(23.8\,\mathrm{MeV}\)
(\texttt{EQ-WH-BOUND-01}). Lean \texttt{RadiationBounds} /
\texttt{EQ-WH-FC-02} further proves the single-mode displaced-oscillator
Franck--Condon class misses the target: \(S^N/N! < 10^{-11}\) for
\textbf{all} \(S\le10^{4}\), \(N\ge300048\) (alias
{\ttfamily harmonic\_\allowbreak{}dos\_\allowbreak{}non\_\allowbreak{}closure}); the lattice
\(\Gamma_{\mathrm{ph}}/\Gamma_{\gamma}\ge10^{11}\) target is not reached
by harmonic DOS or this FC class. This is a \textbf{model-class
non-closure} result---not a proof that radiation suppression is
impossible. Beyond those classes, anharmonic, Umklapp/defect,
multi-mode, and Hagelstein lossy spin-boson mechanisms remain named open
hypotheses. Detector upper limits (\texttt{EQ-RAD-DET-01} DEF mirrors)
rule out only \emph{assumed} branching scenarios; they do not supply a
measured branching ratio ({\ttfamily detector\_\allowbreak{}ul\_\allowbreak{}not\_\allowbreak{}proof}). Cd-ratio /
moderator-null / soft-X∧¬hard-γ PREDs (PRED-08/09/10) are branching
falsifiers---they do \textbf{not} close \texttt{EQ-OPEN-RAD-01}.
\textbf{Falsifier / closure:} a licensed absorption model outside the
bounded classes giving
\(\Gamma_{\mathrm{ph}}/\Gamma_{\gamma}\ge10^{11}\), \textbf{or} measured
\(\mathrm{BR}_{\gamma}\le\) UL with calibrated efficiency + coincidence
provenance.

\hypertarget{a.3-statistical-feasibility-and-rate-closure-q3}{%
\subsubsection{5A.3 Statistical Feasibility and Rate Closure
(Q3)}\label{a.3-statistical-feasibility-and-rate-closure-q3}}

Given the enhancement factors derived in Section 2.8
(\(f_{\text{total}} \sim 10^3\text{-}10^5\)) applied to bare D+D
tunneling probabilities at thermal energies
(\(P_{\text{bare}} \sim 10^{-85}\)):

\[
P_{\text{enhanced}} = P_{\text{bare}} \times f_{\text{total}} \sim 10^{-85} \times 10^5 = 10^{-80}
\]

For a system with \(10^{23}\) atoms cycling at \(10^{12}\) Hz for
\(10^5\) s:

\[
N_{\text{events}} = 10^{23} \times 10^{12} \times 10^5 \times 10^{-80} = 10^{-40}
\]

This predicts zero observable events---many orders of magnitude below
any detection threshold.

\textbf{Resolution pathways:}

\begin{enumerate}
\def\labelenumi{\arabic{enumi}.}
\item
  \textbf{Enhancement factors are vastly underestimated:} The screening
  enhancement alone, if calculated with measured anomalous screening
  energies (\(U_e \sim 300\text{-}800\) eV per Raiola et al.; RR-007
  material-dependent cited ranges), evaluated with the Assenbaum
  screened-cross-section form \(f = \exp(\pi \eta U_e / E)\) at thermal
  energies, is formally enormous. That expansion is only valid for
  \(U_e \ll E\) (Lean/IR machine check: \(U_e < E_{\mathrm{kin}}\),
  \texttt{EQ-ASSEN-DOM-01}); outside the domain the companion software
  still emits a capped sensitivity factor but sets
  {\ttfamily formula\_\allowbreak{}claimed=\allowbreak{}false} --- it does \textbf{not} claim a
  physical Assenbaum enhancement. Converting the screening potential to
  a tunneling enhancement through the proper Gamow integral may yield
  factors exceeding the adopted \(10^3\text{-}10^5\) band, which is why
  reported values are capped there. Do not conflate Assenbaum domain
  validity with the separate RR-007 adiabatic gas-target floor
  (\(U_e \sim 20\text{-}25\) eV).
\item
  \textbf{Collective/coherent multi-body processes:} If LENR proceeds
  through coherent multi-body tunneling (Takahashi's TSC model) or
  condensate transitions (Section 2.6), the bare probability is not
  \(P_{\text{bare}} = e^{-2\pi\eta}\) for a single pair, but a
  collective probability with fundamentally different scaling.
\item
  \textbf{Threshold behavior:} The system may exhibit sharp
  phase-transition-like behavior where below a critical condition the
  rate is essentially zero, but above it the rate jumps discontinuously.
  This is consistent with the observed poor reproducibility---most
  experiments are below threshold.
\end{enumerate}

\textbf{Conclusion on Q3 (\texttt{EQ-OPEN-RATE-01} stays OPEN; Lean
posture mixed):} Rate closure is the most stringent quantitative test of
the framework. Under vacuum-like two-body bookkeeping the lattice rate
budget does not close (\texttt{EQ-RATE-2B-NC-01},
\(\log_{10}N=-40\))---a \textbf{negative determination for that
mechanism form}, not for LENR in general (UI may show
{\ttfamily BUDGET\_\allowbreak{}DOES\_\allowbreak{}NOT\_\allowbreak{}CLOSE} with {\ttfamily established\_\allowbreak{}reference}
tier; the obligation remains OPEN). Under stated engine caps, D--D
trapped states at FCC tetrahedral spacing (or wider) and three-body
capture outside the Widom--Larsen mass-renormalization domain also do
not close (\texttt{EQ-RATE-CAP-NC-01} \(\log_{10}N\le-15\);
\texttt{EQ-RATE-CAP-NC-02} \(\log_{10}N\le-266\); form-DISPROVEN under
caps). For a 25 pm trapped state and for three-body capture at
\(m^*/m_e=4\), engine defaults give \textbf{positive} budgets
(\(\log_{10}N\approx+9.7\) and \(\approx+7.1\)), so caps cannot certify
non-closure there ({\ttfamily EQ-\allowbreak{}RATE-\allowbreak{}CAP-\allowbreak{}GAP-\allowbreak{}01}---a \textbf{proved gap},
not closure); those positives rest entirely on unconstrained inputs (25
pm separation, \(f_{\mathrm{nuc}}=1\), assumed \(m^*/m_e=4\)).
Hydrodynamic collapse enhancement and scenario priors are formally
excluded from counting as nuclear-rate proof (\texttt{EQ-BBL-NR-01}).
WH-ENP-01 uses the same-cluster FGR versus
\(R_{\mathrm{obs}}=P_{\mathrm{W}}/(QV)\) (\texttt{EQ-WH-FLOOR-01});
{\ttfamily when.\allowbreak{}threshold\_\allowbreak{}crossed} is true only when
\(R_{\mathrm{mode}}\ge R_{\mathrm{obs}}\) with \(V\), \(\tau\),
\(P_{\mathrm{occ}}\), and \(f_{\mathrm{nuc}}\) labeled cited vs model
assumption. The v2 closure demand makes the shortfall explicit: for
\(1\,\mathrm{W}\) in a nominal \(1\,\mathrm{cm^{3}}\),
\(R_{\mathrm{obs}}\approx2.6\times10^{17}\,\mathrm{m^{-3}s^{-1}}\); with
the rate engine's own site convention
\(n_{\mathrm{sites}}=0.5\,(n_e/Z)\,\theta\approx6.7\times10^{26}\,\mathrm{m^{-3}}\)
(about \(46\times\) below \(0.5\,n_D\) at the physical Pd atom density;
recorded, not hidden) and
\(\Gamma_{\mathrm{tun}}(25\,\mathrm{pm})\approx6.8\times10^{-37}\), the
required
\(P_{\mathrm{occ}}f_{\mathrm{nuc}}C_{\mathrm{lin}}/\tau\approx5.8\times10^{26}\,\mathrm{s^{-1}}\)
exceeds the engine cap product \(10^{26}\,\mathrm{s^{-1}}\)
(\texttt{EQ-WH-CAP-01}). At every cap the assumed 25 pm trap reaches
\(R_{\max}\approx4.5\times10^{16}\,\mathrm{m^{-3}s^{-1}}\) (factor
\(\approx5.8\) short) and FCC tetrahedral geometry
(\(137.5\,\mathrm{pm}\),
\(\Gamma_{\mathrm{tun}}\approx1.4\times10^{-85}\)) reaches
\(\approx10^{-32}\,\mathrm{m^{-3}s^{-1}}\) (\(\approx49\) decades
short); neither crosses \(R_{\mathrm{obs}}\). The
\((\theta,\tau,f_{\mathrm{nuc}})\) grid peaks at
\(\log_{10}(R_{\mathrm{mode}}/R_{\mathrm{obs}})=-7.4\) and is a scenario
table, not a posterior. Rate-UQ / Oracle B ExpHom identities
(\texttt{EQ-UQ-BF-01}) and the governance enhancement prior upper edge
\(10^{5}\) remain algebraic honesty floors---not laboratory tunneling
measurements or empirical Bayes factors. Overlap DEF
(\texttt{EQ-OVERLAP-01/02}) records ensemble-count bookkeeping; ensemble
product ≠ nuclear rate. \textbf{Until an experimentally constrained
matrix element and state-occupation probability replace the assumed
trapped separation, lifetime, and collective factor, multi-body
(\texttt{trapped\_state}) rate outputs are scenario explorations rather
than predictions.} The software default remains the conservative
\textbf{standard two-body reference baseline} (\texttt{two\_body}):
established free-particle formalism used as a reference, \emph{not} a
validated lattice-LENR mechanism. \textbf{Falsifier / closure:} a named
rate mode giving \(\log_{10}N_{\mathrm{events}}\ge0\) with
\textbf{lab-constrained} \(P_{\mathrm{occ}}\cdot f_{\mathrm{nuc}}\) (or
measured trap separation / matrix element), then governed data +
ledger/matrix PR + sign-off.

\hypertarget{a.4-electron-capture-energy-source-q4}{%
\subsubsection{5A.4 Electron Capture Energy Source
(Q4)}\label{a.4-electron-capture-energy-source-q4}}

The Widom-Larsen mechanism requires electrons with kinetic energy
\(> 0.782\) MeV for \(p + e^- \to n + \nu_e\). Double-layer fields
provide \(\sim 0.1\) eV---a factor of \(\sim 10^7\) deficit.

\textbf{Analysis:}

\begin{enumerate}
\def\labelenumi{\arabic{enumi}.}
\item
  \textbf{Widom-Larsen heavy electron mechanism:} Proposes that
  collective electromagnetic fluctuations at the metal-hydrogen
  interface renormalize the electron effective mass to
  \(m_e^* > 2.53 \, m_e\). This is a very large mass renormalization by
  condensed matter standards (typical electron effective masses in
  metals are \(0.5\text{-}5 \, m_e\)). Whether coherent electromagnetic
  fields can achieve this remains debated. In particular, Tennfors (Eur.
  Phys. J. Plus 128, 15, 2013) estimated, from neutron-scattering data
  on protons in palladium, an electron mass enhancement of less than one
  percent---far below the threshold required---and Ciuchi et
  al.~similarly questioned the predicted rates; Widom and Larsen rebut
  these by arguing the process is collective (many-body) rather than
  two-body. We therefore treat this mechanism as
  \textbf{speculative/contested} and do not rely on it for the paper's
  primary energy-balance argument.
\item
  \textbf{Alternative: D+D priming reactions:} The D+D → ³He + n channel
  (3.27 MeV, athermal neutron) could provide ``seed'' neutrons at low
  rates, which then catalyze Oppenheimer-Phillips transmutation chains
  on heavier lattice nuclei. This avoids the electron capture threshold
  entirely but requires explaining the initial D+D reaction.
\item
  \textbf{Coherent lattice excitation:} If a lattice phonon field
  achieves sufficient coherence (e.g., through the mechanisms in Section
  2.6), the collective energy available at a single site could in
  principle exceed the single-phonon energy by a factor of
  \(\sqrt{N_{\text{coherent}}}\). For \(N \sim 10^8\) coherent phonons,
  this gives \(\sqrt{10^8} \times 50\) meV \(\approx 500\) eV---still
  far below 0.782 MeV.
\end{enumerate}

\textbf{Conclusion on Q4 (\texttt{EQ-OPEN-EC-01} stays OPEN; Lean
posture BOUNDED / form-DISPROVEN):} No currently proposed mechanism
provides a convincing quantitative pathway to produce electrons with
\(E > 0.782\) MeV in condensed matter systems. A \textasciitilde0.1 eV
double-layer source falls short of the 0.782 MeV threshold by
\(\sim7.8\times10^{6}\) ({\ttfamily EQ-\allowbreak{}OPEN-\allowbreak{}EC-\allowbreak{}DEFICIT-\allowbreak{}01}); no
condensed-matter path to the threshold is claimed. Under the stated
per-atom condensed-matter budget (10--100 eV/atom, an assumed AX cap),
Lean {\ttfamily EQ-\allowbreak{}OPEN-\allowbreak{}EC-\allowbreak{}BOUND-\allowbreak{}01} proves no source in that form closes
the threshold; closure would require coherently concentrating at least
7,820 such budgets onto one capture
({\ttfamily sub\_\allowbreak{}minimum\_\allowbreak{}concentration\_\allowbreak{}does\_\allowbreak{}not\_\allowbreak{}close}), which no
mechanism in the framework establishes. This rules out the
\textbf{per-atom form only}; {\ttfamily q4\_\allowbreak{}not\_\allowbreak{}refuted\_\allowbreak{}outright} keeps
the obligation itself open. The D+D priming pathway (option 2) avoids
this problem entirely and is the most physically defensible route to
neutron production. The Widom-Larsen mechanism, while elegant, requires
independent experimental verification of the mass renormalization claim
(UI ``0.000 MeV remaining'' under assumed \(m^*/m_e\) is scenario
bookkeeping, not discharge). Lean \texttt{WeakInteraction} / IR
\texttt{EQ-EC-*} document Fermi identities only. Contested Parkhomov /
detector-gas SPEC curated rows never feed production rates and are not
EC closure. \textbf{Falsifier / closure:} a licensed source outside the
per-atom form delivering \(\ge0.782\,\mathrm{MeV}\) to one capture (or
physical \(m^*/m_e\ge2.53\) within Tennfors domain if Widom--Larsen),
then governed evidence + ledger/matrix PR + sign-off.

\hypertarget{a.5-energy-product-consistency-q5}{%
\subsubsection{5A.5 Energy-Product Consistency
(Q5)}\label{a.5-energy-product-consistency-q5}}

As detailed in Section 5.2.2, the Miles heat-helium correlation
\((2.6 \pm 0.6) \times 10^{11}\) ⁴He/s/W is within a factor of
\textasciitilde2 of the 23.8 MeV/⁴He prediction. Under claim tier
{\ttfamily empirical\_\allowbreak{}lenr\_\allowbreak{}correlation}, that band is the
\textbf{best-studied quantitative EMP correlation} motivating a D+D →
⁴He + lattice-heat \emph{hypothesis}---not discharged proof of the
channel, and not closure of \texttt{EQ-OPEN-YIELD-01} (Lean
{\ttfamily EQ-\allowbreak{}OPEN-\allowbreak{}YIELD-\allowbreak{}BAND-\allowbreak{}01} DEF still OPEN).

\textbf{Conclusion on Q5 (\texttt{EQ-OPEN-YIELD-01} stays OPEN; Lean
posture BOUNDED):} Instantaneous Miles heat--helium correlation
(\textasciitilde2×) is distinct from integrated yield, which remains
\(10^{2}\)--\(10^{4}\times\) low. Lean {\ttfamily EQ-\allowbreak{}OPEN-\allowbreak{}YIELD-\allowbreak{}NC-\allowbreak{}01}
proves that membership in the instantaneous Miles band does \textbf{not}
discharge the integrated-yield obligation
({\ttfamily instantaneous\_\allowbreak{}miles\_\allowbreak{}band\_\allowbreak{}does\_\allowbreak{}not\_\allowbreak{}discharge\_\allowbreak{}q5},
{\ttfamily miles\_\allowbreak{}band\_\allowbreak{}does\_\allowbreak{}not\_\allowbreak{}imply\_\allowbreak{}q5\_\allowbreak{}closed});
\texttt{YieldAccounting} (\texttt{EQ-YIELD-ACCT-01}) mirrors
expected-yield duty-cycle bookkeeping with no observed-yield input.
Integrated yields per W·hr are 100--10,000× lower than expected (Section
5.2.2), indicating either: (a) time-varying reaction rates with duty
cycles, (b) incomplete helium collection, or (c) contributions from
non-helium-producing channels. \textbf{Falsifier / closure:}
time-resolved simultaneous heat + \(^4\mathrm{He}\) with calorimetry and
He-assay efficiency/retention provenance yielding integrated
\(\Delta N_{^4\mathrm{He}}=0\) shortfall (or equivalent licensed
discharge); instantaneous band membership provably does not count.

\begin{center}\rule{0.5\linewidth}{0.5pt}\end{center}

\hypertarget{conclusion}{%
\subsection{6. Conclusion}\label{conclusion}}

This work presents a unified theoretical and mathematical framework that
addresses key LENR phenomena---anomalous heat production, isotopic
shifts, minimal radiation, and neutron generation---through a synthesis
of Rydberg Matter, halo electrons, double-layer QED effects, Casimir
effect, and quantum tunneling. Lattice-bubble interaction physics
(heterogeneous nucleation, bubble--surface coupling, boundary layers,
acoustic forcing) extends modified bubble dynamics to surface-coupled
cavitation. Computational models (Poisson--Schrödinger coupling,
nucleation and surface-reaction integration) organize these mechanisms
under claim tiers; reported excess-heat and transmutation correlations
remain disputed and are not treated as validation of the framework.

\textbf{Future Directions}

\begin{enumerate}
\def\labelenumi{\arabic{enumi}.}
\tightlist
\item
  \textbf{Refinement of Computational Models:} Incorporate many-body
  physics, non-linear effects, and real-time electron--phonon coupling
  under explicit claim tiers (modelled vs aspirational).\\
\item
  \textbf{Independent Falsification Campaigns:} Expand simultaneous
  heat--helium / detector campaigns and PRED falsifiers so OPEN
  obligations can be closed or rejected on evidence---not scaled toward
  commercial viability (energetic/practical freedom remains a permanent
  non-goal).\\
\item
  \textbf{Interdisciplinary Collaboration:} Encourage a synergy of
  nuclear engineers, condensed matter physicists, chemists,
  experimentalists, and computational scientists to resolve outstanding
  LENR challenges without overclaiming nuclear detection.
\end{enumerate}

\hypertarget{mechanized-consistency-lean}{%
\subsubsection{6.1 Mechanized consistency
(Lean)}\label{mechanized-consistency-lean}}

A companion formalization pilot ({\ttfamily formalization/\allowbreak{}lean}, Lean 4
\textbf{without mathlib}; \textbf{no \texttt{sorry}}; exactly
\textbf{one} explicitly labelled bridge axiom
{\ttfamily RealShim.\allowbreak{}expZ\_\allowbreak{}strict\_\allowbreak{}mono}, plus finite
\texttt{native\_decide} checks that trust the Lean compiler)
mechanically checks \textbf{mathematical consistency} of the adopted
classical core: Rydberg \(n=10\) identity, Casimir check-geometry band
(\(\lvert F\rvert\approx 1.3\,\mathrm{mN}\)), Gamow penetration
anti-monotonicity, Assenbaum domain predicate
(\(U_e < E_{\mathrm{kin}}\)), enhancement composition with cap
\(10^{5}\) and the illegal two-engine product (\texttt{EQ-COMP-03}), the
Miles \(Q\)-value conversion identity, phonon order \(N=Q/\hbar\omega\),
two-body rate-budget does-not-close (\texttt{EQ-RATE-2B-NC-01}),
scenario cap bounds / gaps ({\ttfamily EQ-\allowbreak{}RATE-\allowbreak{}CAP-\allowbreak{}NC-\allowbreak{}01/\allowbreak{}02},
{\ttfamily EQ-\allowbreak{}RATE-\allowbreak{}CAP-\allowbreak{}GAP-\allowbreak{}01}), MeV remainder caps
(\texttt{EQ-WH-MEV-CAP-01}), FC/harmonic class radiation bounds
(\texttt{EQ-WH-FC-02}), EC CM-budget family
({\ttfamily EQ-\allowbreak{}OPEN-\allowbreak{}EC-\allowbreak{}BOUND-\allowbreak{}01}), yield instantaneous≠integrated
({\ttfamily EQ-\allowbreak{}OPEN-\allowbreak{}YIELD-\allowbreak{}NC-\allowbreak{}01}), hydro≠nuclear-rate
(\texttt{EQ-BBL-NR-01}), n-body / Overlap geometry vs trap identities
(\texttt{EQ-OVERLAP-01/02}; ensemble ≠ rate), collapse \textbf{clamp /
UQ / RP-domain / CR-cap} bookkeeping (\texttt{CollapseBounds} /
\texttt{EQ-BBL-*}---not a cavitation-dynamics theorem), WeakInteraction
Fermi \(G_F\)/σ/Boltzmann bounds (\texttt{EQ-EC-*}; ≠
\texttt{EQ-OPEN-EC-01} closure), Decision ExpHom / \texttt{EQ-UQ-BF-01}
algebraic BF floors (not empirical Bayes), and aspirational
Poisson--Schrödinger continuum DEF
(\texttt{EQ-PS-01}/\texttt{EQ-PS-TOL-01}; never public-beta FEM). An
equation intermediate representation (\texttt{formalization/ir/}) links
paper statements, Python literals, Wolfram numeric checks, and Lean
targets. Lean non-closure / bound language (WH-ENP-01, OpenObligations,
OpenQuestionBounds) remains the honesty floor: \texttt{lenr\_proved}
stays false and every \texttt{EQ-OPEN-*} stays OPEN until a later
evidence-backed closure PR. Theorem names that mention ``closure'' prove
non-closure or contrapositives under stated caps/classes; none
discharges an obligation.

\textbf{Multi-format paper pack (v0.3.7-beta).} Academic archival copies
of this manuscript and the Lean section→theorem crosswalk ship as:

\begin{itemize}
\tightlist
\item
  Markdown / plain text / LaTeX / PDF:
  {\ttfamily LENR\_\allowbreak{}SIM\_\allowbreak{}APP\_\allowbreak{}v2/\allowbreak{}docs/\allowbreak{}v0.\allowbreak{}3.\allowbreak{}7-\allowbreak{}beta/\allowbreak{}paper/\allowbreak{}} (rebuild via
  \texttt{build\_pdf.sh})
\item
  Lean crosswalk module:
  {\ttfamily formalization/\allowbreak{}lean/\allowbreak{}LenrFormal/\allowbreak{}PaperCrosswalk.\allowbreak{}lean} (mirrored
  in the paper pack)
\item
  Equation LaTeX stubs:
  {\ttfamily LENR\_\allowbreak{}SIM\_\allowbreak{}APP\_\allowbreak{}v2/\allowbreak{}docs/\allowbreak{}v0.\allowbreak{}3.\allowbreak{}7-\allowbreak{}beta/\allowbreak{}latex/\allowbreak{}}
\item
  NL↔MD↔JSON↔LaTeX↔Lean matrix:
  {\ttfamily LENR\_\allowbreak{}SIM\_\allowbreak{}APP\_\allowbreak{}v2/\allowbreak{}docs/\allowbreak{}v0.\allowbreak{}3.\allowbreak{}7-\allowbreak{}beta/\allowbreak{}PROOF\_\allowbreak{}MATRIX.\allowbreak{}md}
\end{itemize}

\textbf{Non-claims.} These theorems do \textbf{not} prove that LENR
occurs, that contested calorimetry is experimentally true, or that the
open challenges of §5A (MeV gap, radiation suppression, rate closure,
electron-capture threshold, integrated yield) are solved. Speculative
mechanisms remain claim-tiered. Status ledger:
{\ttfamily docs/\allowbreak{}FORMALIZATION/\allowbreak{}FORMAL\_\allowbreak{}PROOF\_\allowbreak{}STATUS.\allowbreak{}md}.

\textbf{Would-close-if lemmas (still OPEN).}
{\ttfamily LenrFormal.\allowbreak{}OpenQuestionBounds} records closure \emph{predicates}
and integer-scaled implications only: Q1 closes if the energy remainder
is identically zero (Hagelstein/Dicke/BEC remain named hypotheses); Q2
would close if the lattice \(\Gamma_{\mathrm{ph}}/\Gamma_{\gamma}\)
bookkeeping ratio meets \(10^{11}\); Q3 occupancy/collective factors are
Nat inputs to the existing {\ttfamily RateBudget.\allowbreak{}would\_\allowbreak{}close\_\allowbreak{}if} product
(two-body \(\log_{10}N=-40\) remains a non-closure theorem); Q4 proves a
millielectronvolt \textbf{deficit} of \(\sim 10^{7}\) between a stated
0.1 eV double-layer and the 0.782 MeV capture threshold
({\ttfamily EQ-\allowbreak{}OPEN-\allowbreak{}EC-\allowbreak{}DEFICIT-\allowbreak{}01}) and a CM-budget form-DISPROVEN family
({\ttfamily EQ-\allowbreak{}OPEN-\allowbreak{}EC-\allowbreak{}BOUND-\allowbreak{}01}) without discharging
\texttt{EQ-OPEN-EC-01}; Q5 distinguishes an instantaneous Miles EMP band
vs integrated 100--\(10^{4}\) remainder ({\ttfamily EQ-\allowbreak{}OPEN-\allowbreak{}YIELD-\allowbreak{}BAND-\allowbreak{}01})
and proves instantaneous band membership does not discharge integrated
yield ({\ttfamily EQ-\allowbreak{}OPEN-\allowbreak{}YIELD-\allowbreak{}NC-\allowbreak{}01}) without discharging parent
obligation \texttt{EQ-OPEN-YIELD-01}.
{\ttfamily LenrFormal.\allowbreak{}WorkingHypothesis} (WH-ENP-01) proves a coupled
remainder vector and \textbf{non-closure} (\texttt{EQ-WH-NONCLOSE-01})
without flipping the OPEN table.
{\ttfamily LenrFormal.\allowbreak{}WorkingHypothesisBounds} (WH-ENP-01 v2) adds
integer-scaled bounds: the Poisson step ratio \(S/(N+1)<1\) for
\(N\ge S\) with \(\log_{10}P_N\le-7\times10^{5}<-11\) on the
\(S\le10^{4}\) grid (\texttt{EQ-WH-FC-01}; conservative cap on the
\texttt{lgamma} value); the general partition
\(E_{\mathrm{derived}}+(Q-E_{\mathrm{derived}})=Q\) with remainder
\(=0\iff E_{\mathrm{derived}}=Q\) (\texttt{EQ-WH-PART-02});
\(\lfloor n_D V_{\mathrm{dicke}}\rfloor=0\), hence \(\neg(n\le0)\) for
\(n\ge2\) (\texttt{EQ-WH-DICKE-N-01}); the capped product
\(P_{\mathrm{occ}}f_{\mathrm{nuc}}C_{\mathrm{lin}}\le1\cdot1\cdot10^{6}\)
so a demand above the cap is unmet (\texttt{EQ-WH-CAP-01}); and the
conjunction \(\neg Q_1\wedge\neg Q_2\wedge\neg Q_3\)
(\texttt{EQ-WH-NONCLOSE-02}). The Franck--Condon model is a labelled
model assumption, not a nature theorem. There is no
{\ttfamily theorem\ lenr\_\allowbreak{}occurs}.

\begin{center}\rule{0.5\linewidth}{0.5pt}\end{center}

\hypertarget{mathematical-approaches-and-their-relevance-to-lenr}{%
\subsection{7. Mathematical Approaches and Their Relevance to
LENR}\label{mathematical-approaches-and-their-relevance-to-lenr}}

While the preceding sections introduced key theoretical constructs, LENR
also relies heavily on rigorous mathematical methods for both modeling
and empirical validation. Below is a concise overview:

\begin{enumerate}
\def\labelenumi{\arabic{enumi}.}
\tightlist
\item
  \textbf{Quantum Mechanics and Tunneling}
\end{enumerate}

\begin{itemize}
\tightlist
\item
  \textbf{Schrödinger Equation:} \[
   -\frac{\hbar^2}{2m}\nabla^2\psi + V(r)\psi = E\psi,
   \] central for calculating wavefunction overlaps, tunneling rates,
  and delocalized-interface-electron interactions. Numerical integration
  (finite difference, finite element) helps solve complex potential
  landscapes.
\end{itemize}

\begin{enumerate}
\def\labelenumi{\arabic{enumi}.}
\setcounter{enumi}{1}
\tightlist
\item
  \textbf{Thermodynamics and Energy Conservation}
\end{enumerate}

\begin{itemize}
\tightlist
\item
  \textbf{Excess Heat Models:} \[
   Q_{\text{output}} = Q_{\text{input}} + Q_{\text{nuclear}}, \qquad Q_{\text{excess}} = Q_{\text{output}} - Q_{\text{input}},
   \] ensuring total energy balance in calorimetric setups.
\item
  \textbf{Partition Functions:} \[
   Z = \sum_i e^{-E_i/k_B T},
   \] used for analyzing statistical behavior of lattice atoms, phonons,
  and quasi-particles.
\end{itemize}

\begin{enumerate}
\def\labelenumi{\arabic{enumi}.}
\setcounter{enumi}{2}
\tightlist
\item
  \textbf{Field Theory and Electrodynamics}
\end{enumerate}

\begin{itemize}
\tightlist
\item
  \textbf{Extended Maxwell's Equations:} \[
   \nabla\cdot\mathbf{E} = \frac{\rho}{\epsilon_0}, \qquad \nabla\times\mathbf{B} = \mu_0\mathbf{J} + \frac{1}{c^2}\frac{\partial\mathbf{E}}{\partial t},
   \] crucial for modeling plasma-like oscillations and double-layer
  fields in electrochemical or gaseous LENR environments.
\end{itemize}

\begin{enumerate}
\def\labelenumi{\arabic{enumi}.}
\setcounter{enumi}{3}
\tightlist
\item
  \textbf{Nuclear Reaction Cross-Sections}
\end{enumerate}

\begin{itemize}
\tightlist
\item
  \textbf{Reaction Rate Expression:} \[
   R_{\text{target}} = n\,\sigma_{\text{nuc}}\,v, \qquad R_V = n_1 n_2\langle\sigma_{\text{nuc}} v\rangle
   \] with \(R_{\text{target}}\) a per-target rate (s\(^{-1}\)) and
  \(R_V\) the volumetric rate (m\(^{-3}\) s\(^{-1}\)). For identical
  deuterons use \(R_V=\tfrac12 n_D^2\langle\sigma v\rangle\). LENR
  enhancements arise from small effective barriers (via Casimir, high
  local fields, or RM).
\end{itemize}

\begin{enumerate}
\def\labelenumi{\arabic{enumi}.}
\setcounter{enumi}{4}
\tightlist
\item
  \textbf{Lattice Dynamics}
\end{enumerate}

\begin{itemize}
\tightlist
\item
  \textbf{Harmonic Oscillators:} \[
   H = \sum_i\left(\frac{p_i^2}{2m} + \frac{1}{2}k x_i^2\right),
   \] used to approximate phonons interacting with embedded deuterons or
  protons.
\item
  \textbf{Cluster Fusion Models:} Tetrahedral Symmetric Condensate (TSC)
  → advanced group-theoretic tools describe multi-body fusion in metals
  (Takahashi's approach; claim tier speculative).
\end{itemize}

\begin{enumerate}
\def\labelenumi{\arabic{enumi}.}
\setcounter{enumi}{5}
\tightlist
\item
  \textbf{Casimir and Vacuum Fluctuations}
\end{enumerate}

\begin{itemize}
\tightlist
\item
  \textbf{Casimir Force:} \[
   F_{\text{Casimir}} = -\frac{\pi^2 \hbar c}{240}\frac{A}{d^4},
   \] leading to an attractive potential that can enhance
  electron--nucleus proximity (correlation ≠ nuclear proof).
\end{itemize}

\begin{enumerate}
\def\labelenumi{\arabic{enumi}.}
\setcounter{enumi}{6}
\item
  \textbf{Machine Learning and Statistical Analysis}

  \begin{itemize}
  \tightlist
  \item
    \textbf{Pattern Recognition:} Clustering and neural networks to
    identify parameter regimes that maximize excess heat or
    transmutation yields.\\
  \item
    \textbf{Monte Carlo Simulations:} Stochastic exploration of reaction
    pathways, guiding experimental setups.\\
  \end{itemize}
\item
  \textbf{Numerical Validation}

  \begin{itemize}
  \tightlist
  \item
    \textbf{Molecular Dynamics (MD):} Tracks atom-by-atom collisions,
    temperature spikes, and bubble collapse phenomena.\\
  \item
    \textbf{Finite Element Analysis (FEA):} Solves boundary value
    problems for electric fields, thermal gradients, and mechanical
    stress in reactor components.
  \end{itemize}
\end{enumerate}

Through this collection of mathematical and computational tools,
researchers systematically link theoretical LENR constructs with
real-world experimental observations, thereby advancing the rigor,
reproducibility, and credibility of the field.

\begin{center}\rule{0.5\linewidth}{0.5pt}\end{center}

\hypertarget{glossary-of-key-mathematical-symbols}{%
\subsection{8. Glossary of Key Mathematical
Symbols}\label{glossary-of-key-mathematical-symbols}}

Below is a reference list defining the primary symbols and terms used in
this paper:

\begin{itemize}
\tightlist
\item
  \(\boldsymbol{\nabla^2}\): The Laplacian operator,
  \(\nabla^2 = \frac{\partial^2}{\partial x^2} + \frac{\partial^2}{\partial y^2} + \frac{\partial^2}{\partial z^2}\).
\item
  \(\boldsymbol{\phi}\): Electric potential \([\text{V}]\).
\item
  \(\boldsymbol{\rho}\): Charge density \([\text{C/m}^3]\).
\item
  \(\boldsymbol{\epsilon_0}\): Vacuum permittivity,
  \(\approx 8.854 \times 10^{-12}\,\text{F/m}\).
\item
  \(\boldsymbol{\psi}\): Wavefunction describing the quantum state of a
  particle.
\item
  \(\boldsymbol{V(\mathbf{r})}\): Potential energy function
  \([\text{J}]\) or \([\text{eV}]\).
\item
  \(\boldsymbol{E}\): (1) Energy eigenvalue in the Schrödinger equation;
  (2) Electric field strength \([\text{V/m}]\), depending on context.
\item
  \(\boldsymbol{\hbar}\): Reduced Planck's constant,
  \(\hbar = h/2\pi \approx 1.054 \times 10^{-34}\,\text{J}\cdot\text{s}\).
\item
  \(\boldsymbol{c}\): Speed of light in vacuum,
  \(c \approx 3 \times 10^8\,\text{m/s}\).
\item
  \(\boldsymbol{m}\): Mass (electron, proton, or other particle),
  context-specific.
\item
  \(\boldsymbol{k_B}\): Boltzmann's constant,
  \(k_B \approx 1.381 \times 10^{-23}\,\text{J/K}\).
\item
  \(\boldsymbol{a_0}\): Bohr radius, \(a_0 \approx 0.0529\,\text{nm}\).
\item
  \(\boldsymbol{n}\): (1) Principal quantum number; (2) Number density
  \([\text{m}^{-3}]\), depending on context.
\item
  \(\boldsymbol{n_1, n_2}\): Number densities of two (possibly
  identical) reacting species \([\text{m}^{-3}]\); use
  \(\tfrac{1}{2}n^2\) rather than \(n_1 n_2\) when the species are
  identical (§5.1.1).
\item
  \(\boldsymbol{N}\): Number density of target nuclei (OP process, §2.7)
  \([\text{m}^{-3}]\).
\item
  \(\boldsymbol{\Phi_n}\): Neutron flux \([\text{m}^{-2}\text{s}^{-1}]\)
  (§2.7).
\item
  \(\boldsymbol{a}\): Characteristic distance scale (e.g., extended
  orbital radius).
\item
  \(\boldsymbol{r_0}\): Nuclear core or reference radius.
\item
  \(\boldsymbol{P_{\text{overlap}}}\): Overlap probability between
  interface electron and nuclear wavefunctions.
\item
  \(\boldsymbol{p}\): Proton, or momentum operator (\(\hat{p}\))
  depending on context.
\item
  \(\boldsymbol{e^-}\): Electron.
\item
  \(\boldsymbol{\nu_e}\): Electron neutrino.
\item
  \(\boldsymbol{F_{\text{Casimir}}}\): Casimir force.
\item
  \(\boldsymbol{A}\): Area (e.g., Casimir effect) \([\text{m}^2]\).
\item
  \(\boldsymbol{d}\): Separation distance (e.g., Casimir or
  double-layer) \([\text{m}]\).
\item
  \(\boldsymbol{T}\): Temperature \([\text{K}]\).
\item
  \(\boldsymbol{T_c}\): Critical temperature for Bose--Einstein
  condensation.
\item
  \(\boldsymbol{\zeta(3/2)}\): Riemann zeta function at \(3/2\),
  \(\zeta(3/2) \approx 2.612\).
\item
  \(\boldsymbol{\Delta V}\): Potential difference \([\text{V}]\).
\item
  \(\boldsymbol{\mu}\): (1) Dynamic viscosity
  \([\text{Pa}\cdot\text{s}]\); (2) reduced mass \([\text{kg}]\),
  depending on context.
\item
  \(\boldsymbol{\sigma}\): (1) Generic cross-section \([\text{m}^2]\);
  (2) surface tension \([\text{N/m}]\) in Rayleigh--Plesset. Context
  clarifies usage.
\item
  \(\boldsymbol{\sigma_0}\): Base/reference cross-section
  \([\text{m}^2]\).
\item
  \(\boldsymbol{\sigma_{\text{nuc}}}\): Nuclear (fusion) cross-section
  used in the generic rate expression
  \(R_{\text{rate}} = n\sigma_{\text{nuc}}v\) (§7) \([\text{m}^2]\).
\item
  \(\boldsymbol{\sigma_{\text{OP}}}\): Oppenheimer--Phillips capture
  cross-section (§2.7) \([\text{m}^2]\).
\item
  \(\boldsymbol{\sigma_{\text{surf}}}\): Bubble surface tension in the
  Rayleigh--Plesset equation (§3.2) \([\text{N/m}]\).
\item
  \(\boldsymbol{\langle\sigma v\rangle}\): Thermally-averaged
  cross-section--velocity product (§5.1.1) \([\text{m}^3/\text{s}]\).
\item
  \(\boldsymbol{R}\): (1) Reaction rate \textbf{per unit volume},
  \([\text{reactions}\cdot\text{m}^{-3}\text{s}^{-1}]\) --- see
  §2.7/§5.1.1/§7 for three distinct regimes of this quantity; (2) Bubble
  radius in the Rayleigh--Plesset equation \([\text{m}]\).
\item
  \(\boldsymbol{P_{\text{tunnel}}}\): Quantum tunneling (Gamow)
  probability, dimensionless.
\item
  \(\boldsymbol{P_{\text{threshold}}}\): Probability of exceeding
  reaction threshold conditions, dimensionless (§5.1.1).
\item
  \(\boldsymbol{f_{\text{enhancement}}}\): Cumulative
  tunneling-enhancement factor, dimensionless;
  \(f_{\text{enhancement}} = f_{\text{screening}} \times f_{\text{field}} \times f_{\text{coherence}} \times f_{\text{phonon}}\)
  (§2.8, §5.1.1).
\item
  \(\boldsymbol{Q_{\text{nuclear}}, Q_{\text{input}}, Q_{\text{excess}}}\):
  Energy quantities in calorimetry \([\text{J}]\).
\item
  \(\boldsymbol{\zeta}\): Generic variable for partition functions or
  the Riemann zeta function, as context requires.
\end{itemize}

\textbf{Formalization symbol map (IL-015).} In the Lean/IR ledger
({\ttfamily docs/\allowbreak{}FORMALIZATION/\allowbreak{}CLAIM\_\allowbreak{}TAXONOMY.\allowbreak{}md}), overloaded symbols are
disambiguated as \texttt{E\_eigen} / \texttt{E\_field} /
\texttt{E\_kin}, \texttt{n\_prin} / \texttt{n\_dens}, and
\texttt{R\_rate} / \texttt{R\_bubble} / \texttt{R\_overlap\_count}.

\begin{center}\rule{0.5\linewidth}{0.5pt}\end{center}

\hypertarget{references}{%
\subsection{References}\label{references}}

Citations below support the mathematical modeling in this paper.
Preference is given to sources tied to \textbf{replicable experimental
reports} (calorimetry, helium/isotope assays, screening measurements)
and to \textbf{established mechanisms} used in the equations (tunneling,
Casimir, Oppenheimer--Phillips, phonon--nuclear coupling). Speculative
mechanisms remain claim-tiered in the body text even when cited.

\hypertarget{subtitle-homage}{%
\subsubsection{Subtitle homage}\label{subtitle-homage}}

\textbf{{[}vaswani2017{]}} Vaswani, A., Shazeer, N., Parmar, N.,
Uszkoreit, J., Jones, L., Gomez, A. N., Kaiser, Ł., \& Polosukhin, I.
(2017). Attention Is All You Need. \emph{Advances in Neural Information
Processing Systems} 30.\\
\emph{Methodological metaphor only---not a LENR physics source.} The
subtitle ``Coupling is all you need'' deliberately echoes this work's
framing of coupled multi-agent representations.\\
\textbf{Link:} \url{https://arxiv.org/abs/1706.03762} ·
\href{https://papers.nips.cc/paper_files/paper/2017/hash/3f5ee243547dee91fbd053c1c4a845aa-Abstract.html}{NeurIPS
2017}

\hypertarget{primary-sources-for-mathematical-modeling-replicable-evidence-established-mechanisms}{%
\subsubsection{Primary sources for mathematical modeling (replicable
evidence \& established
mechanisms)}\label{primary-sources-for-mathematical-modeling-replicable-evidence-established-mechanisms}}

Drawn from the project's evidence-base bibliography
({\ttfamily Unified\_\allowbreak{}LENR\_\allowbreak{}Theory\_\allowbreak{}D.\allowbreak{}Acs.\allowbreak{}docx.\allowbreak{}md}) and from DOI-backed
screening / heat--helium constraints used by the companion simulation.

\begin{enumerate}
\def\labelenumi{\arabic{enumi}.}
\item
  Fleischmann, M., \& Pons, S. (1989). Electrochemically induced nuclear
  fusion of deuterium. \emph{Journal of Electroanalytical Chemistry},
  261(2), 301--308.\\
  \textbf{Link:} \url{https://doi.org/10.1016/0022-0728(89)80006-3}
\item
  McKubre, M. C. H., et al.~(1994). Excess power observations in
  electrochemical studies of the D/Pd system. \emph{Journal of
  Electroanalytical Chemistry}, 368, 55--66.\\
  \textbf{Link:} \url{https://doi.org/10.1016/0022-0728(94)87011-6}
\item
  Storms, E. (2010). A critical analysis of the LENR phenomenon.
  \emph{Naturwissenschaften}, 97(10), 861--881.\\
  \textbf{Link:} \url{https://doi.org/10.1007/s00114-010-0711-x}
\item
  Arata, Y., \& Zhang, Y.-C. (2002). Observation of anomalous heat
  release and helium-4 production from highly deuterated palladium fine
  particles. \emph{Proceedings of the Japan Academy, Series B}, 78(7),
  57--61.\\
  \textbf{Link:} \url{https://doi.org/10.2183/pjab.78.57}
\item
  Mizuno, T. (2007). Anomalous energy generation during treatment of
  nickel with hydrogen gas. \emph{Journal of Nuclear Science and
  Technology}, 44(2), 331--338.\\
  \textbf{Link:} \url{https://doi.org/10.1080/18811248.2007.9711293}
\item
  Mizuno, T., \& Rothwell, J. (2019). Increased excess heat from
  palladium deposited on nickel. \emph{Journal of Condensed Matter
  Nuclear Science}, 29, 1--30.\\
  \textbf{Link:} \url{http://www.iscmns.org/CMNS/JCMNS-Vol29.pdf}
\item
  Miles, M. H., et al.~(1993). Correlation of excess power and helium
  production. \emph{Journal of Electroanalytical Chemistry}, 346,
  99--117.\\
  \textbf{Link:} \url{https://doi.org/10.1016/0022-0728(93)85700-H}
\item
  Steinetz, B. M., et al.~(2020). Novel nuclear reactions observed in
  bremsstrahlung-irradiated deuterated metals. \emph{Physical Review C},
  101, 044610.\\
  \textbf{Link:} \url{https://doi.org/10.1103/PhysRevC.101.044610}
\item
  Pines, V., et al.~(2020). Nuclear fusion reactions in deuterated
  metals. \emph{Physical Review C}, 101, 044609.\\
  \textbf{Link:} \url{https://doi.org/10.1103/PhysRevC.101.044609}
\item
  Raiola, F., et al.~(2002). Electron screening in d(d,p)t for
  deuterated metals. \emph{European Physical Journal A}, 13, 377--382.\\
  \textbf{Link:} \url{https://doi.org/10.1007/s10050-002-8766-5}
\item
  Huke, A., et al.~(2008). Enhancement of deuteron-fusion reactions in
  metals. \emph{Physical Review C}, 78, 015803.\\
  \textbf{Link:} \url{https://doi.org/10.1103/PhysRevC.78.015803}
\item
  Czerski, K., et al.~(2023). Electron screening in palladium.
  \emph{Physics Letters B}, 845, 137750.\\
  \textbf{Link:} \url{https://doi.org/10.1016/j.physletb.2023.137750}
\item
  Fralick, G., Decher, R., \& Rousar, D. (2003). Anomalous heat effects
  in metal hydrides. NASA/TM---2003-212795.\\
  \textbf{Link:} \url{https://ntrs.nasa.gov/citations/20040152112}
\item
  Iwamura, Y., Itoh, T., \& Takahashi, T. (2024). Anomalous heat
  generation in nano-composite hydrides. \emph{Japanese Journal of
  Applied Physics}, 63(3), 037001.\\
  \textbf{Link:} \url{https://doi.org/10.35848/1347-4065/ad2622}
\item
  Bellucci, S., Cardone, F., \& Pistella, F. (2023). Review of
  experiments reporting non-conventional phenomena in nuclear matter.
  \emph{Symmetry}, 15(8), 1507.\\
  \textbf{Link:} \url{https://doi.org/10.3390/sym15081507}
\item
  Oppenheimer, J. R., \& Phillips, M. (1935). Note on the transmutation
  function for deuterons. \emph{Physical Review}, 48(7), 500--502.\\
  \textbf{Link:} \url{https://doi.org/10.1103/PhysRev.48.500}
\item
  Bordag, M., Mohideen, U., \& Mostepanenko, V. M. (2001). New
  developments in the Casimir effect. \emph{Physics Reports}, 353(1--3),
  1--205.\\
  \textbf{Link:} \url{https://doi.org/10.1016/S0370-1573(01)00015-1}
\item
  Hagelstein, P. L., \& Chaudhary, I. U. (2015). Coherent energy
  exchange between lattice and nuclear degrees of freedom. \emph{Journal
  of Condensed Matter Nuclear Science}, 15, 1--24.\\
  \textbf{Link:} \url{http://www.iscmns.org/CMNS/JCMNS-Vol15.pdf}
\item
  Takahashi, A. (2004). Tetrahedral symmetric condensate fusion model
  and multi-body fusion in deuterated metal. \emph{Journal of Condensed
  Matter Nuclear Science}, 1, 4--17.\\
  \textbf{Link:} \url{http://www.iscmns.org/CMNS/JCMNS-Vol1.pdf}
\item
  Takahashi, A. (2016). Fundamental of rate theory for CMNS.
  \emph{Journal of Condensed Matter Nuclear Science}, 19, 298--315.
  (Basis for the software \texttt{trapped\_state} rate mode.)
\item
  Holmlid, L. (2015). Rydberg Matter and its detection: A review of
  experimental techniques. \emph{Journal of Physics: Conference Series},
  635, 012002.\\
  \textbf{Link:} \url{https://doi.org/10.1088/1742-6596/635/1/012002}
\item
  Widom, A., \& Larsen, L. (2006). Ultra low-momentum neutron catalyzed
  nuclear reactions on metallic hydride surfaces. \emph{The European
  Physical Journal C}, 46(1), 107--111.\\
  \textbf{Link:} \url{https://doi.org/10.1140/epjc/s2006-02479-8}
\item
  Preparata, G. (1995). \emph{QED Coherence in Matter}. World
  Scientific.\\
  \textbf{Link:}
  \url{https://www.worldscientific.com/worldscibooks/10.1142/2739}
\item
  Schwinger, J. (1990). Cold fusion: A hypothesis. \emph{Zeitschrift für
  Naturforschung A}, 45, 756--758.\\
  \textbf{Link:}
  \url{https://www.lenr-canr.org/acrobat/SchwingerJcoldfusion.pdf}
\item
  Kim, Y. E., \& Zubarev, A. L. (2001). Theory of Bose--Einstein
  condensation mechanism for deuteron-induced nuclear reactions in
  micro/nano-scale metal grains and particles. \emph{Physical Review A},
  64, 013603.\\
  \textbf{Link:} \url{https://doi.org/10.1103/PhysRevA.64.013603}
\item
  Taleyarkhan, R. P., et al.~(2002). Evidence for nuclear emissions
  during acoustic cavitation. \emph{Science}, 295(5561), 1868--1873.\\
  \textbf{Link:} \url{https://doi.org/10.1126/science.1067589}
  \emph{Honesty:} Contested / disputed nuclear-emission claim; retained
  as acoustic-cavitation literature anchor only---not laboratory LENR
  proof. See also \texttt{pyyhtia2026} (electrolytic CR-39 cavitation
  confounder).
\item
  Menlove, H. O., et al.~(1989). Neutron emission from titanium exposed
  to deuterium gas at high pressure. Los Alamos National Laboratory
  Report.\\
  \textbf{Link:} \url{https://www.osti.gov/biblio/5039188}
\item
  Bockris, J. O'M., Lin, G. H., \& Bush, B. F. (1991). Tritium and
  neutron generation in electrolysis of heavy water. \emph{Proceedings
  of the Electrochemical Society} / LENR-CANR archive.\\
  \textbf{Link:}
  \url{https://lenr-canr.org/acrobat/BockrisJdoestritiu.pdf}
\item
  Wolf, K. L., et al.~(1990). Tritium production in deuterated
  palladium. Texas A\&M / related reports.\\
  \textbf{Link:}
  \url{https://link.springer.com/article/10.1007/BF02627575}
\item
  Celani, F., et al.~(2014). Experimental results on sub-micro
  structured Pd wires and thin films. \emph{Journal of Condensed Matter
  Nuclear Science}, 13, 56--68.\\
  \textbf{Link:} \url{http://www.iscmns.org/CMNS/JCMNS-Vol13.pdf}
\item
  Bakranov, N., Didenko, V., \& Nagel, D. (2024). Nanomaterials
  engineering for enhanced LENR: A comprehensive review. \emph{Frontiers
  in Materials}, 11, 1500487.\\
  \textbf{Link:} \url{https://doi.org/10.3389/fmats.2024.1500487}
\item
  Levi, G., et al.~(2013). Indications of anomalous heat energy
  production in a reactor device containing hydrogen-loaded nickel
  powder. arXiv:1305.3913.\\
  \textbf{Link:} \url{https://arxiv.org/abs/1305.3913}
\item
  Del Giudice, E., Doglia, S., Milani, M., \& Vitiello, G. (1985).
  Electrodynamical coherence in biological systems. \emph{Physical
  Review Letters}, 54(7), 622--625.\\
  \textbf{Link:} \url{https://doi.org/10.1103/PhysRevLett.54.622}
\item
  Safajuei, F., Safajuei, S., \& Hajimohammadi, M. (2023). Production of
  ultra-low momentum neutrons on metallic hydride surfaces by heavy
  electrons. \emph{Iranian Journal of Science and Technology,
  Transactions A: Science}, 47, 995--1002.\\
  \textbf{Link:} \url{https://doi.org/10.1007/s40995-023-01437-3}
\end{enumerate}

\hypertarget{supplementary-keys-software-uq-traceability}{%
\subsubsection{Supplementary keys (software / UQ
traceability)}\label{supplementary-keys-software-uq-traceability}}

These \texttt{AuthorYear} keys resolve claims in the companion app and
{\ttfamily docs/\allowbreak{}REVIEW\_\allowbreak{}2026/\allowbreak{}TRACEABILITY\_\allowbreak{}MATRIX.\allowbreak{}md}:

\begin{itemize}
\tightlist
\item
  \textbf{{[}steinetz2020exp{]}} → entry 8 above\\
\item
  \textbf{{[}pines2020theory{]}} → entry 9 above\\
\item
  \textbf{{[}raiola2002{]}} → entry 10 above\\
\item
  \textbf{{[}huke2008{]}} → entry 11 above\\
\item
  \textbf{{[}czerski2023{]}} → entry 12 above\\
\item
  \textbf{{[}miles1993{]}} → entry 7 above\\
\item
  \textbf{{[}takahashi2016{]}} → entry 20 above\\
\item
  \textbf{{[}hagelstein2025{]}} Hagelstein, P. L. (2025). Relativistic
  phonon--nuclear coupling / lossy spin-boson reviews (MeV-gap
  channel).\\
\item
  \textbf{{[}screeningreview2024{]}} Electron screening in laboratory
  nuclear reactions. \emph{Particles}, 7(3), 50. DOI:
  \href{https://doi.org/10.3390/particles7030050}{10.3390/particles7030050}\\
\item
  \textbf{{[}storms2015{]}} Storms, E. (2015). Efficient production of
  energy from the LENR effect. \emph{Current Science}, 108(4),
  574--577.\\
\item
  \textbf{{[}tennfors2013{]}} Tennfors, E. (2013). On the assumptions
  for an effective electron mass. \emph{European Physical Journal Plus},
  128, 15. DOI:
  \href{https://doi.org/10.1140/epjp/i2013-13015-3}{10.1140/epjp/i2013-13015-3}\\
\item
  \textbf{{[}mulnix2026{]}} Mulnix, A. (2026). Structural theory of
  sonoluminescence via Klein--Gordon / Rayleigh--Plesset harmonic
  locking. Zenodo. DOI:
  \href{https://doi.org/10.5281/zenodo.19079074}{10.5281/zenodo.19079074}\\
\item
  \textbf{{[}pyyhtia2026{]}} Pyyhtiä, K., et al.~(2026). Investigation
  of CR-39 damaging mechanisms in electrochemical environments.
  \emph{Journal of Electroanalytical Chemistry}, 1013, 120104. DOI:
  \href{https://doi.org/10.1016/j.jelechem.2026.120104}{10.1016/j.jelechem.2026.120104}.
  OA experimental support that HER nanobubble cavitation damages CR-39
  and confounds pit=nuclear interpretations; does \textbf{not} prove
  LENR.\\
\item
  \textbf{{[}nguyen2026{]}} Nguyen, T., et al.~(2026). Quantum shell
  structure of proton--neutron pairing. \emph{Nature}. DOI:
  \href{https://doi.org/10.1038/s41586-026-10616-2}{10.1038/s41586-026-10616-2}\\
\item
  \textbf{{[}floquetwkb2026{]}} Floquet--WKB analysis of lattice
  confinement fusion under a driven barrier (2026). DOI:
  \href{https://doi.org/10.14293/PR2199.002776.v1}{10.14293/PR2199.002776.v1}\\
\item
  \textbf{{[}parkhomov2018{]}} Parkhomov, A. G. (2018). Nuclear
  reactions in LENR: review of experimental results (LENR Academy
  curated subset).
\end{itemize}

\begin{center}\rule{0.5\linewidth}{0.5pt}\end{center}

\end{document}
